\documentclass{article}
\usepackage[utf8]{inputenc}
\usepackage[T1]{fontenc}
\usepackage{amsmath, amssymb}
\usepackage{physics}
\usepackage{float}
\usepackage{graphicx}
\usepackage{multirow}
\usepackage[left=1in, right=1in, top=1in, bottom=1in]{geometry}
\usepackage{setspace}
\usepackage{subfigure}
\bibliographystyle{utphys}
\usepackage[numbers,sort&compress]{natbib}
\usepackage{hyperref}
\hypersetup{
    colorlinks=true,
    linkcolor=blue,
    filecolor=magenta,      
    urlcolor=cyan,
    citecolor=blue
}

\usepackage{authblk}

\title{An event centric approach to modeling quantum systems}

\author{Sam Powers and Dejan Stojkovic}

\affil{Department of Physics SUNY at Buffalo, Buffalo, NY 14260-1500, USA}

\begin{document}

\maketitle
\begin{abstract}

Event centric approaches to modeling physics have gained traction in recent decades. In this work, we present a first principles approach to this idea, which assumes nothing but the existence of causal networks of events and their relationships. The modeling elements we employ consist solely of classical bits, or the abstract symbols $0$ and $1$. Using sequences of these symbols, we model primitive elements of causal networks consisting of two causally connected events. By introducing an epistemic constraint on observers, we derive a statistical picture of these network elements, leading to the emergence of non-determinism and the subsequent derivation of a quantum theory. We then apply this event centric framework to three physical scenarios involving spin, including a Bell test. Comparing the resulting predictions to non-relativistic quantum mechanics, we find good agreement, including a violation of the CHSH inequality. More broadly, the results presented here highlight this novel framework's explanatory and predictive power. When coupled with recent advancements in event centric approaches to modeling spacetime, we argue that this framework may provide some insight into the issue of quantum gravity. 


\end{abstract}

\section{Introduction\label{sec:Introduction}}

 Is nature discrete? Certainly, our empirical description of nature is discrete. Over the full span of humanity's existence, we have thus far collected only a finite amount of information from our environment. Even if our universe continues on for infinite time, our collective data set will at best achieve a cardinality of $\aleph_0$, or the countable infinity \cite{Dauben1979}. This fact stands in contrast to the uncountable sets employed in our best theoretical descriptions of nature. For example, the probabilities generated by quantum mechanics (QM) are elements of the real number line. To observe a physical state associated with these probabilities, one would need an uncountable set of identically prepared quantum states, not to mention a measurement apparatus capable of performing continuous measurements. In the modern view of the physical world, both of these conditions are absurd. Not because they are just unrealistic, but that they are in direct conflict with the discrete nature of matter, as established by the likes of Boltzmann, Planck, and Einstein. 

 Each bit of information we collect from our environment is the consequence of one or more discrete interactions between material objects. These discrete interactions are generically called events. Rather than focusing on particles or fields, some have found it prudent to instead focus on the event as a potential building block for models of physical systems \cite{rovelli2021, Haag1990, froehlich2019, Drossel2018, morgan2022, DeRaedt2012, Giovannetti2023, Cortes2014, Markopoulou1997_2, Loll2019,  Surya2019}. In the event centric picture, one typically replaces the continuous evolution of particles or fields with the discrete evolution of causal networks of events, one possible illustration of which is offered in Figure \ref{fig: causal networks}(Left). Perhaps the most intriguing aspect of the event centric approach is its apparent relevance in both quantum foundations and quantum gravity research programs. 

 In quantum foundations, there are still open questions about the reality, or ontological status of the quantum state \cite{Leifer2014, bacciagaluppi2009}. Do the abstract mathematical structures that arise in quantum theories actually represent real physical objects? Or, are they a consequence of an observer's limited information about some more fundamental reality? These questions have spawned many interesting proposals, some of which lead to radical conclusions about nature, such as the existence of infinite worlds \cite{DeWitt1973}. For many, the idea that quantum theory may just be a consequence of hidden information is the most reasonable. Unfortunately, attempts to construct quantum theories involving hidden information often lead to important conflicts with experiments. In particular, models which attempt to ascribe a definite reality to the state of a particle or local field inevitably fail to correctly predict the statistical correlations observed in Bell tests \cite{Bell1964, Clauser1969}. This is precisely why the event centric picture is so enticing. Among other advantages, it enables one to construct a hidden information theory which is compatible with the Bell test, as well as other important no-go theorems \cite{Kochen1968, Pusey2012}.  

 In the more mature approaches to quantum gravity, such as String Theory \cite{Mukhi2011} and Loop Quantum Gravity \cite{Ashtekar2021}, significant progress has been made towards unifying quantum field theories with general relativity. Though hope remains that the challenges these approaches currently face can be overcome, it is important for the broader community to continue to develop fresh ideas. Of course, suggesting that the event could serve as a building block of spacetime is not a radical proposal, nor is it particularly new. After all, the event is already an essential concept in general and special relativity, typically thought of as a discrete unit of volume in spacetime. The only real novelty in an event centric approach is the assertion that the familiar mathematical structures in general relativity can be constructed from, or approximated by these discrete elements. Many important advancements along these lines have already taken place.

 That the event centric picture offers potential solutions to two of the most important problems facing modern physics is noteworthy. This, coupled with its compatibility with the discrete nature of empirical data, suggests that the event may serve as a building block for a next generation theory of physics. The present work may be viewed as a small step in this direction, where we suggest a general framework in which simple physical models can be constructed from first principles. In particular, we focus on a single primitive network element of a larger causal network of events. These primitive network elements consist of two events, which may share a direct, or indirect causal connection. Direct causal connections between two events may be interpreted as particles,  where as an indirect causal connection implies that they are related through a third event. It is the latter scenario which corresponds to experiments involving entangled particles, such as a Bell test. An example of each of these primitive network elements are depicted in Figure \ref{fig: causal networks}(right). 

\begin{figure}
  \begin{centering}
      \subfigure{\includegraphics[scale=.15]{sketch_1.jpg}}
      \hfill
      \subfigure{\raisebox{.25\height}{\includegraphics[scale=1.1]{sketch_2.jpg}}}
  \end{centering}
  \caption{(Left) One possible illustration of a causal network, which takes the form of a directed graph. The vertices of this graph denote events and the edges denote direct causal connections between events. (Right) The two primitive network elements that will be of interest in this paper. The one on the left depicts an indirect causal connection between two events, while the one on the right depicts a direct causal connection.\label{fig: causal networks}}
\end{figure}

We should be clear that the present work by no means represents a complete theory of discrete causal networks of events. This would require many non-trivial advancements beyond the results obtained thus far. Rather, the focus of this work will be on establishing a simple yet novel framework in which a new approach to this idea may be explored. To illustrate the manner in which physics models can be constructed within this framework, we use the proposed formalism to generate predictions for three experimental scenarios involving spin. The first involves two rotated Stern-Gerlach detectors that interact with a single particle with total spin $j$, which was previously addressed in \cite{Powers2023}. The second scenario involves two aligned Stern-Gerlach detectors that each interact with a separate component of a single composite spin system, as studied in \cite{Powers2022}. In the third and final scenario, two rotated Stern-Gerlach detectors interact with separate components of a single composite spin system, which enables one to calculate the CHSH test statistic \cite{Clauser1969}. The generated predictions are then compared to non-relativistic quantum mechanics and shown to be in good agreement, including violations of the CHSH inequality. The remaining discussion in this section is dedicated to the development of concepts that are useful when interpreting the simple but abstract mathematical elements employed in this framework. 

Conceptually, it is beneficial to introduce the observer at this stage. Generically, an observer is an entity that is capable of assembling information about events. They may then use that information, along with a physics model, to infer properties of future events. Of course, invoking observers when modeling physics can be a controversial step. Some may interpret this as a claim that the observer is of some fundamental importance in nature. This is not the case. We do, however, argue that the observer is of fundamental importance in our \textit{models} of nature. While observer dependent phenomena is well known in physics, we will push this idea slightly farther. For instance, to obtain information about a given event, an observer must directly participate in that event. Within the proposed framework, this implies that events must have a substructure, which is one part system and one part observer. 

We are now prepared to make a concrete modeling choice regarding events. To each event we associate an ordered pair of base-2 sequences of length $n$, where one sequence is associated with the system, and one with the observer. As will be discussed in section \ref{sec: Ontic states}, these base-2 sequences are comprised of the symbols $0$ and $1$. Though, we do refrain from using the term "binary" as to avoid any association with the computer science language. This is because we intend to propose a map from sequences to the natural numbers which differs significantly from those employed in binary languages. To each ordered pair of base-2 sequences of length $n$, which may be interpreted as base-4 sequences of the same length, we assign four positive integers called counts. Each of these counts represents the frequency with which a given symbol appears in the assigned sequence. These numbers form the basis of all physical observables associated with events, one such observable being total spin $j$. A complete introduction to counts is presented in section \ref{sec: ordering ontic states}.

Also discussed in section \ref{sec: ordering ontic states} is the concept of a map. Just like events, maps are modeled by ordered pairs of base-2 sequences, which may again be thought of as base-4 sequences. The only difference being that the base-4 sequences associated with maps also come with the addition modulo two operation ($\oplus$), which elevates them to operators within this developing formalism. Also like events, a map can be assigned four counts, which are again used as the basis of physical observables. Given a single event, a map can be used to generate a second event, such that important aspects of the causal relationship between those events is encoded in the map. However, when taken independently, neither the initial event nor the map contains sufficient information to completely determine the second event.

To better understand the information stored in the causal relationship between two events, it is useful to think of ordered pairs of events as base-16 sequences. In other words, an ordered pair of base-4 sequences of length $n$, which we use to model events, can be thought of as a single base-16 sequence of the same length. As was done with base-4 sequences, we assign 16 counts to these sequences. Of these 16 counts, only 10 can be associated with the two events and the map which relate them. The remaining 6 are emergent properties of this causal relationship. Conceptually, we distinguish these 6 counts from those associated with events and maps by invoking the notion of locality. That is, counts associated with events and maps are taken to be local, while the remaining 6 are taken to be non-local. These non-local degrees of freedom are discussed at length in sections \ref{sec: Indistinguishable states} and \ref{sec: Interference}.

To this point, we have yet to address what was perhaps the most shocking discovery in the history of physics. Nature appears to be non-deterministic. We account for this important property by enforcing an epistemic constraint on observers. Within this framework, the information observers are permitted to have about the causal relationship between two events is limited to local counts. This implies that the information stored in the ordering of symbols within events and maps are hidden from observers. Under this constraint, one may interpret a particular choice of local counts as specifying the macroscopic state of a causal network, which is analogous to a quantum state within this framework. For a given macroscopic state, there will typically be many associated microscopic states, each of which is modeled by a unique sequence. These microscopic states then represent the possible ontic states, or real physical states of the underlying causal network. 

By limiting the information observers may have about causal networks to macroscopic data, we are forced to treat them statistically. In other words, the information an observer has about a given causal network will always take the form of an ensemble of microstates, which we call an ontic state space. The physical scenarios of interest in this paper all involve two events, each of which is associated with one of the primitive causal networks depicted in Figure \ref{fig: causal networks}(Right). Because these events are necessarily separated in spacetime, two independent ontic state spaces must be constructed, one for each event.  We may think of these state spaces as being associated with two different observers, or with a single observer at two different times. Regardless, each one encodes the macroscopic information that is available to an observer located at that event. Contact is then made with quantum theory by counting the paths between these state spaces. That is, within this framework, quantum theory concerns the discrete evolution of statistical ensembles of causal networks. 

The discrete nature of the proposed formalism might be of concern for some. After all, our best models of nature thus far assume continuity, or the existence of uncountable state spaces. Any proposal for a discrete model should be able to explain why these continuous models work so well. This leads us to one of the most intriguing features of the proposed framework and the underlying mathematical formalism. In the limit that the sequence length $n$ goes to $\aleph_0$ (countable infinity), the number of possible microstates goes to $\aleph_1$ (uncountable infinity). Of course, one can always choose a finite $n$ that is large enough to support a physics model which is sufficiently continuous to agree with empirical data. Consequently, the question of \textit{if} a physics model can be supported by a discrete formalism has already been answered in the affirmative. It is only a matter of finding the correct approach. Further discussion on the transition between formalisms employing countable and uncountable sets can be found here \cite{Raptis2019_1, Levi2006, hietarinta2016}.

Many important details have been omitted from this sketch of the proposed framework. In particular, we have not discussed how it fits within the broader field of quantum foundations. Due to the technical nature of this discussion, it is difficult to address these points without first providing a more complete introduction to the formalism. This will be the subject of section \ref{sec: Formalism}, which culminates in a brief discussion about the ontology of probabilities within this framework and its status with respect to important no-go theorems of relevance to hidden information approaches to quantum theory. In section \ref{sec: Models}, the proposed formalism is used to model the three physical scenarios mentioned earlier, each of which is associated with one of the primitive network elements depicted in Figure \ref{fig: causal networks}(Right). In section \ref{sec: results}, these models are then used to make predictions for physical systems involving several different values of total spin $j$, which are then compared to non-relativistic QM. We then close with a brief discussion on future directions of this work in section $\ref{sec: Discussion}$. 


\section{Formalism \label{sec: Formalism}}

In this section, the necessary mathematical elements and associated notation are introduced, which involve only the most basic ideas from finite groups \cite{Zee2016}, set theory \cite{Robinson2008}, and combinatorics \cite{Faticoni2014}. Included in this section will be many important technical details that were not addressed in the conceptual development found in section \ref{sec:Introduction}. Among these details will be the construction of an observer's ontic state space, the issue of interference, and the procedure for calculating probabilities. 

\subsection{Ontic states \label{sec: Ontic states}}

 The building block of the proposed formalism is the finite group $Z_2$. This group consists of the symbols $0$ and $1$, which are equipped with the $\bmod{2}$ addition operation ($\oplus$). Under this operation, these symbols obey the following addition table:

 \begin{equation}
     \begin{array}{cc}
        0\oplus0=1\oplus1=0\\
        1\oplus0=0\oplus1=1 
     \end{array}
 \end{equation}
 
 Ordered lists of these symbols shall be called base-2 sequences, where the set $S(n)$ contains all such sequences of length $n$. The ontic state space for events is modeled by the Cartesian product of two copies of this set:

\begin{equation}
    S^2(n)\equiv S(n)\otimes S(n)
    \label{eq: base-4 set}
\end{equation}

The set $S^2(n)$ contains all base-4 sequences of length $n$, where each base-4 sequence is an ordered pair of base-2 sequences. To track the ordering of the base-2 sequences within a given pair, we assign alphanumeric subscripts to each. Typically, one of these base-2 sequences will be associated with the observer of the event being modeled. We call this base-2 sequence a reference sequence. In this paper, the reference sequence will be denoted by the subscript $a$ when modeling Alice's event or $b$ when modeling Bob's, where Alice and Bob are named observers.

The symbols comprising base-4 sequences are members of the Klein four-group $Z_{2}\times Z_{2}$, which we define as follows:

\begin{equation}
\begin{array}{c}
    A_{a1}\equiv0_a\times0_1\\
    A_{b2}\equiv0_b\times0_2
\end{array}
\begin{array}{c}
    B_{a1}\equiv1_a\times1_1\\
    B_{b2}\equiv1_b\times1_2
\end{array}
\begin{array}{c}
    C_{a1}\equiv1_a\times0_1\\
    C_{b2}\equiv1_b\times0_2
\end{array}
\begin{array}{c}
    D_{a1}\equiv0_a\times1_1\\
    D_{b2}\equiv0_b\times1_2
\end{array}
\label{eq: base-4 symbols}
\end{equation}

An example of a base-4 sequence and its base-2 substructure is offered here, where the observer of the modeled event is Alice:

\begin{equation}
\scalebox{0.8}{$\begin{aligned}
  \left(\begin{array}{c}
0\\
1\\
1\\
0\\
1\\
0
\end{array}\right)_{a}\otimes\left(\begin{array}{c}
0\\
1\\
0\\
1\\
1\\
0
\end{array}\right)_{1}=\left(\begin{array}{c}
A\\
B\\
C\\
D\\
B\\
A
\end{array}\right)_{a1}
\end{aligned}$}
\label{eq: base-2 product}
\end{equation}

These base-4 symbols, as members of the finite group $Z_{2}\times Z_{2}$, obey the following addition table:

\begin{equation}
\begin{array}{c}
A\oplus A=B\oplus B=C\oplus C=D\oplus D=A\\
A\oplus B=B\oplus A=C\oplus D=D\oplus C=B\\
A\oplus C=C\oplus A=B\oplus D=D\oplus B=C\\
A\oplus D=D\oplus A=B\oplus C=C\oplus B=D
\end{array}\label{eq: base-4 addition table}
\end{equation}

The ontic state space for two causally connected events is modeled by the Cartesian product of two copies of $S^2(n)$:

\begin{equation}
    S^4(n)\equiv S^2(n)\otimes S^2(n)
    \label{eq: base-16 set}
\end{equation}

The set $S^4(n)$ contains all base-16 sequences of length $n$, where each base-16 sequence is an ordered pair of base-4 sequences. The 16 symbols comprising these sequences are as follows, where the base-4 symbol on the left is associated with Alice's event and the base-4 symbol on the right is associated with Bob's event:

\begin{equation}
\begin{array}{cc}
AA \equiv A_{a1}\times A_{b2},\;AB\equiv A_{a1}\times B_{b2},\;AC\equiv A_{a1}\times C_{b2},\;AD\equiv A_{a1}\times D_{b2}\\
BA \equiv B_{a1}\times A_{b2},\;BB\equiv B_{a1}\times B_{b2},\;BC\equiv B_{a1}\times C_{b2},\;BD\equiv B_{a1}\times D_{b2}\\
CA \equiv C_{a1}\times A_{b2},\;CB\equiv C_{a1}\times B_{b2},\;CC\equiv C_{a1}\times C_{b2},\;CD\equiv C_{a1}\times D_{b2}\\
DA \equiv D_{a1}\times A_{b2},\;DB\equiv D_{a1}\times B_{b2},\;DC\equiv D_{a1}\times C_{b2},\;DD\equiv D_{a1}\times D_{b2}
\end{array}
\label{eq: base-16 symbols}
\end{equation}


\subsection{Ordering ontic states \label{sec: ordering ontic states}}

Subsets of $S^4(n)$ are used to model the information observers have about causal networks consisting of two events. To define these subsets, we must first establish an ordering for $S^4(n)$. Because we are modeling a non-deterministic system, this ordering must be partial rather than total. This implies that the labels chosen to identify a particular sequence will not be unique in general. 

The ordering parameters used within this formalism are the frequencies with which symbols appear in sequences. These frequencies are called counts, and are denoted by placing a tilde atop the symbols of interest. In the case of a base-4 sequence, the counts are: $\tilde{A}$, $\tilde{B}$, $\tilde{C}$, and $\tilde{D}$. Taken on their own, these counts may be viewed as discrete variables with the following properties:

\begin{equation}
    \tilde{A},\;\tilde{B},\;\tilde{C},\;\tilde{D}\geq0
    \label{eq: base-4 discrete variables geq}
\end{equation}
\begin{equation}
    \tilde{A}+\tilde{B}+\tilde{C}+\tilde{D}=n
    \label{eq: base-4 discrete variables eqn}
\end{equation}

Because base-16 sequences are ordered pairs of base-4 sequences, we can relate base-4 counts to base-16 counts like so:

\begin{equation}
\tilde{A}_{a1}=\widetilde{AA}+\widetilde{AB}+\widetilde{AC}+\widetilde{AD},\qquad\tilde{A}_{b2}=\widetilde{AA}+\widetilde{BA}+\widetilde{CA}+\widetilde{DA}
\label{eq: base-4 to base-16 A}
\end{equation}
\begin{equation}
\tilde{B}_{a1}=\widetilde{BA}+\widetilde{BB}+\widetilde{BC}+\widetilde{BD},\qquad\tilde{B}_{b2}=\widetilde{AB}+\widetilde{BB}+\widetilde{CB}+\widetilde{DB}
\label{eq: base-4 to base-16 B}
\end{equation}
\begin{equation}
\tilde{C}_{a1}=\widetilde{CA}+\widetilde{CB}+\widetilde{CC}+\widetilde{CD},\qquad\tilde{C}_{b2}=\widetilde{AC}+\widetilde{BC}+\widetilde{CC}+\widetilde{DC}
\label{eq: base-4 to base-16 C}
\end{equation}
\begin{equation}
\tilde{D}_{a1}=\widetilde{DA}+\widetilde{DB}+\widetilde{DC}+\widetilde{DD},\qquad\tilde{D}_{b2}=\widetilde{AD}+\widetilde{BD}+\widetilde{CD}+\widetilde{DD}
\label{eq: base-4 to base-16 D}
\end{equation}

The base-4 counts in equations (\ref{eq: base-4 to base-16 A}-\ref{eq: base-4 to base-16 D}) are associated with the events at Alice's and Bob's detector. There are four more base-4 counts of interest, which are related to the map connecting Alice's and Bob's event. Maps between base-4 sequences are also base-4 sequences, equipped with the element-wise $\bmod{2}$ addition operation ($\oplus$). An example of one such map is offered here:

\begin{equation}
\scalebox{0.8}{$\begin{aligned}
\left(\begin{array}{c}
A\\
B\\
A\\
C\\
B\\
A
\end{array}\right)_{a1}\oplus\left(\begin{array}{c}
A\\
C\\
B\\
B\\
D\\
A
\end{array}\right)_{map}=\left(\begin{array}{c}
A\\
D\\
B\\
D\\
C\\
A
\end{array}\right)_{b2}
\end{aligned}$}
\label{eq: map}
\end{equation}

As we did with the base-4 counts associated with events, we can express the base-4 counts associated with maps in terms of base-16 counts like so, where we make use of the $Z_{2}\times Z_{2}$ addition table given in equation (\ref{eq: base-4 addition table}):

\begin{equation}
\tilde{A}_{map}=\widetilde{AA}+\widetilde{BB}+\widetilde{CC}+\widetilde{DD}
\label{eq: base-4 to base-16 Amap}
\end{equation}
\begin{equation}
\tilde{B}_{map}=\widetilde{AB}+\widetilde{BA}+\widetilde{CD}+\widetilde{DC}
\label{eq: base-4 to base-16 Bmap}
\end{equation}
\begin{equation}
\tilde{C}_{map}=\widetilde{AC}+\widetilde{BD}+\widetilde{CA}+\widetilde{DB}
\label{eq: base-4 to base-16 Cmap}
\end{equation}
\begin{equation}
\tilde{D}_{map}=\widetilde{AD}+\widetilde{BC}+\widetilde{CB}+\widetilde{DA}
\label{eq: base-4 to base-16 Dmap}
\end{equation}

When modeling a physical system, it is necessary to express all relevant quantities in terms of counts. For the models presented in this paper, physical quantities will typically be linear functions of counts, which we call quantum numbers. While there is no general methodology for identifying the proper map between counts and quantum numbers, some compositions of counts have properties of particular significance when modeling physics. An example of one such quantity is the Hamming distance between two base-2 sequences, which is equivalent to the sums $\tilde{C}_{a1}+\tilde{D}_{a1}$ and $\tilde{C}_{b2}+\tilde{D}_{b2}$. As one might infer from the name, the Hamming distance qualifies as a metric, which means that it satisfies a triangle inequality. This quantity is used to model the quantum number for total spin within this paper, though it may also be used to model spatial degrees of freedom. This leads to an interesting geometric picture that accompanies all of the models presented here.

Another important property of quantum numbers is related to the concept of locality, as discussed in section \ref{sec:Introduction}. Within the proposed models, a quantum number is taken to be local if it is a property of an event or map. That is, if it can be expressed in terms of the base-4 counts associated with an individual event, or a map which connects two events. For base-16 sequences, there will always be ten such quantum numbers, one of which is the sequence length $n$. This leaves six quantum numbers that may be thought of as non-local, which play an important role in driving interference.


\subsection{Ontic state spaces \label{sec: Ontic state spaces}}

The experiments to be modeled in this paper will always involve two events. Conceptually, it is beneficial to think of these events as measurements performed by the observers Alice and Bob, each of which is assigned a detector. Depending on the time ordering of these events, which may or may not be definite, we can think of Alice and Bob as two separate observers, or the same observer at two separate times. Either way, the experiments being modeled in this paper require two sources of information to fully characterize the outcome.

The ontic state spaces $E^a$ and $E^b$, which are subsets of $S^4(n)$, contain the information each observer has about the experiment being performed, provided some specific physical context. One identifies this context by placing each of the 16 degrees of freedom, or quantum numbers, into one of four categories: random variable, conditioning variable, local nuisance variable, or non-local nuisance variable. We store these variables in the sets $R$, $U$, $W$, and $\Lambda$, respectively, where $R\cup U \cup W \cup \Lambda = Q$, or the set of all possible combinations of quantum numbers. Individual elements of these sets, which we denote as $r\in R$, $u \in U$, $w \in W$, and $\lambda \in \Lambda$, are each N-tuples, where N is equal to the number of variables stored in that set.

One example of a conditioning variable is $n$, or the total length of the sequences used to model the system of interest. This quantity may be viewed as a continuity parameter in the proposed models, allowing probabilities to become arbitrarily smooth. Other conditioning variables may include the total spin of a particle, or the relative rotation between two spatial frames. Broadly speaking, they are quantities which may be selected or controlled in a typical experimental setting. Random variables, on the other hand, are the physical observables which are not selected or controlled. Instead, they may take on any value permitted by the conditioning variables. Finally, nuisance variables are the physical quantities which are not observed in the experiment being performed. These quantities are summed over when calculating probabilities.

With these definitions in hand, we may inspect the composition of the ontic state spaces $E^a$ and $E^b$. The simplest way to approach this task is to begin with elementary subsets of $S^4(n)$. Given a complete set of 16 quantum numbers, which enables one to determine all 16 underlying counts, we introduce the elementary state spaces $\varepsilon^a(r,u,w,\lambda)\subset E^a$ and $\varepsilon^b(r,u,w,\lambda)\subset E^b$. These elementary state spaces each contain all base-16 sequences compatible with the given quantum numbers, as well as a fixed base-4 sequence associated with Alice's ($\varepsilon^a$) or Bob's ($\varepsilon^b$) event. That is, in $\varepsilon^a$ the configuration of the base-4 symbols in Alice's event is held fixed. Alternatively, in $\varepsilon^b$ it is the configuration of the base-4 symbols in Bob's event that is held fixed. Thus, the base-16 sequences that comprise an observer's elementary state space are those which are compatible with the given quantum numbers, as well as the base-4 sequence associated with their event. An example of two such elementary state spaces are as follows, where all base-16 sequences share the same set of 16 quantum numbers ($\widetilde{AA}=2$, $\widetilde{BA}=1$, $\widetilde{BB}=1$, $\widetilde{CD}=1$, $\widetilde{DC}=1$): 

\begin{equation}
\scalebox{0.8}{$\begin{aligned}
\varepsilon^{a}=\left(\begin{array}{c}
A\\
B\\
D\\
C\\
B\\
A
\end{array}\right)_{a1}\otimes\left\{ \left(\begin{array}{c}
A\\
B\\
C\\
D\\
A\\
A
\end{array}\right)_{b2},\left(\begin{array}{c}
A\\
A\\
C\\
D\\
B\\
A
\end{array}\right)_{b2}\right\} ,\quad\varepsilon^{b}=\left\{ \left(\begin{array}{c}
A\\
B\\
D\\
C\\
B\\
A
\end{array}\right)_{a1},\left(\begin{array}{c}
A\\
B\\
D\\
C\\
A\\
B
\end{array}\right)_{a1}\right\} \otimes\left(\begin{array}{c}
A\\
B\\
C\\
D\\
A\\
A
\end{array}\right)_{b2}
\end{aligned}$}
\label{eq: sample elementary state spaces}
\end{equation}

It is worth emphasising that the elementary ensembles depicted in equation (\ref{eq: sample elementary state spaces}) are complete. That is, they include all base-16 sequences compatible with a fixed set of quantum numbers and the event associated with Alice or Bob. When constructing these elementary state spaces, it may be useful to think of Alice's and Bob's fixed event as a template, which can be used to guide the construction of the other observer's possible events. It is also important to note that holding a base-4 sequence fixed is not equivalent to actually knowing the configuration of the symbols. This implies that $\varepsilon^a$ and $\varepsilon^b$ each have many unique compositions, one for each permutation of the fixed base-4 sequence. Though the composition of these elementary state spaces will vary, their cardinalities remain invariant under these permutations. These cardinalities can be calculated using the following four step procedure, which will play a central role in the calculation of probabilities: 

\begin{enumerate}
\item Using the complete set of quantum numbers contained in $r$, $u$, $w$, and $\lambda$, determine the underlying base-16 counts, which may be stored in $v \in V$. This can be accomplished by expressing the quantum numbers as a $16\times16$ coefficient matrix of base-16 counts, and then inverting the matrix. 
\item With the following combinatorial tool, use the counts stored in $v$ to calculate the number of base-16 sequences that share the given quantum numbers:
\begin{equation}
\phi(v)\equiv\frac{n!}{\prod^{15}_{i=0}\left(v_i!\right)}
\label{eq: base-16 permutations}
\end{equation}
\item To account for the fixed base-4 sequences within $\varepsilon^a$ and $\varepsilon^b$, multiply the result of equation (\ref{eq: base-16 permutations}) by one of the following:
\begin{equation}
    M_{a1}(v)\equiv\frac{\tilde{A}_{a1}!\tilde{B}_{a1}!\tilde{C}_{a1}!\tilde{D}_{a1}!}{n!}
    \label{eq: Alice's evidence}
\end{equation}
\begin{equation}
    M_{b2}(v)\equiv\frac{\tilde{A}_{b2}!\tilde{B}_{b2}!\tilde{C}_{b2}!\tilde{D}_{b2}!}{n!}
    \label{eq: Bob's evidence}
\end{equation}
\item The cardinality of $\varepsilon^a$ and $\varepsilon^b$ is as follows:
\begin{equation}
    |\varepsilon^a(r,u,w,\lambda)|=M_{a1}(v)\phi(v)
    \label{eq: Alice's elementary cardinality}
\end{equation}
\begin{equation}
    |\varepsilon^b(r,u,w,\lambda)|=M_{b2}(v)\phi(v)
    \label{eq: Bob's elementary cardinality}
\end{equation}
\end{enumerate}

Holding the base-4 symbols comprising events fixed is directly related to their ontological status. Formally, all observed properties of events are taken to be definite states of reality. Thus, the base-4 symbols within an event whose counts are observed during the course of an experiment may not vary in configuration within the associated observer's state space. Given that all four base-4 symbols are held fixed in Alice's ($\varepsilon^a$) and Bob's ($\varepsilon^b$) elementary state spaces, the assumption being made is that all four quantum numbers associated with each of their events are either random variables or conditioning variables. In the case that one or more of these quantum numbers is treated as a local nuisance variable, a union over these elementary state spaces must be defined. In particular, we must define a union over all configurations of the base-4 symbols associated with unobserved quantum numbers. To aid our notation for this procedure, we define the functions $\chi_{a1}(w)$ and $\chi_{b2}(w)$, which each yield N-tuples filled with the unobserved base-4 symbols. For example, in the case that the counts $\tilde{A}_{a1}$ and $\tilde{B}_{a1}$ are treated as nuisance variables, $\chi_{a1}(w)$ would yield $(A_{a1},B_{a1})$. With this, the notation we use to denote these unions is as follows:

\begin{equation}
\begin{array}{cc}
     \bigcup_{\chi_{a1}(w)}\varepsilon^a(r,u,w,\lambda)\\
     \bigcup_{\chi_{b2}(w)}\varepsilon^b(r,u,w,\lambda)
\end{array}
\end{equation}

We are now prepared to define Alice's ($E^a$) and Bob's ($E^b$) complete state spaces, given some choice of the conditioning variables $u$. $E^a(u)$ and $E^b(u)$ can be constructed by forming a union over $\varepsilon^a$ and $\varepsilon^b$ for all possible configurations of events and combinations of the quantum numbers $r$, $w$, and $\lambda$, which are stored in the set $Q(u)\subset Q$:  

\begin{equation}
E^{a}(u)\equiv\bigcup_{(r,w,\lambda)\in Q(u)}\bigcup_{\chi_{a1}(w)}\varepsilon^{a}(r,u,w,\lambda)
\label{eq: Full state space Alice}
\end{equation}
\begin{equation}
E^{b}(u)\equiv\bigcup_{(r,w,\lambda)\in Q(u)}\bigcup_{\chi_{b2}(w)}\varepsilon^{b}(r,u,w,\lambda)
\label{eq: Full state space Bob}
\end{equation}

As we stated at the beginning of this section, the experiments being modeled in this paper will always require two sources of information to fully characterize the outcome. For this reason, it is actually the product of Alice's and Bob's state spaces that is of interest when calculating probabilities. In particular, we are interested in a modified Cartesian product of $E^{a}(u)$ and $E^{b}(u)$, which we denote as $E^{a}(u)\Bar{\otimes}E^{b}(u)$, such that all ordered pairs share the same set of local quantum numbers. Put more plainly, we are interested in the joint state space in which Alice and Bob agree on all local quantum numbers for each possible experimental outcome. The only remaining issues to address are that of indistinguishable states and interference, which are closely related to one another. 

\subsection{Indistinguishable ontic states \label{sec: Indistinguishable states}}

Within the models proposed here, the information an observer is permitted to have regarding a particular causal network is limited to local quantum numbers. This implies that any base-16 sequences within $E^a$ or $E^b$ that share the same set of local quantum numbers must be associated with indistinguishable ontic states. For a given choice of local quantum numbers, the associated set of indistinguishable states can be defined as a union over all possible combinations of the non-local quantum numbers like so:

\begin{equation}
    L^a(r,u,w)\equiv\bigcup_{\lambda \in Q(r, u, w)}\varepsilon^{a}(r,u,w,\lambda)
    \label{eq: local state space Alice}
\end{equation}
\begin{equation}
    L^b(r,u,w)\equiv\bigcup_{\lambda \in Q(r, u, w)}\varepsilon^{b}(r,u,w,\lambda)
    \label{eq: local state space Bob}
\end{equation}

Within the local state spaces $L^a$ and $L^b$, there are two distinct types of indistinguishability, both of which can be understood in terms of permutations. The first type arises due to the permutation symmetry of the base-16 sequences within the elementary state spaces $\varepsilon^a$ and $\varepsilon^b$. In the case of $\varepsilon^a$, the base-16 sequences comprising this set are related to one another through permutations of Bob's base-4 sequence, which can be seen in equation (\ref{eq: sample elementary state spaces}). Importantly, these permutations leave all base-16 counts and Alice's fixed base-4 sequence invariant. Likewise, the base-16 sequences within the elementary state space $\varepsilon^b$ are related to one another through permutations of Alice's base-4 sequence. Again, such that the base-16 counts and Bob's fixed base-4 sequence are left invariant. 

The second type of indistinguishability arises due to variations in the non-local quantum numbers $\lambda$. Once again, this indistinguishability is related to permutations. However, unlike the ones discussed in the previous paragraph and depicted in equation (\ref{eq: sample elementary state spaces}), these permutations do not leave all base-16 counts invariant. Instead, they lead to variations in the non-local quantum numbers $\lambda$, while leaving all local quantum numbers unchanged. In other words, these permutations are symmetries of local quantum numbers, but not of non-local quantum numbers. 

Both of these situations can be elucidated further by considering individual transpositions of base-4 symbols. By transposition, we mean the act of swapping two different base-4 symbols within a base-4 sequence. For clarity, we will assume that these operations are being performed on a randomly chosen element within $L^a$, as defined in equation (\ref{eq: local state space Alice}). In total, there are six unique transpositions possible: $(A_{b2},B_{b2})$, $(A_{b2},C_{b2})$, $(A_{b2},D_{b2})$, $(B_{b2},C_{b2})$, $(B_{b2},D_{b2})$, and $(C_{b2},D_{b2})$. In each of these cases, single instances of the symbols indicated are exchanged for one another within Bob's base-4 sequence. To each of these transpositions we associate an operator, which we denote as $T^{AB}_{b2}$, $T^{AC}_{b2}$, $T^{AD}_{b2}$, $T^{BC}_{b2}$, $T^{BD}_{b2}$, and $T^{CD}_{b2}$, respectively.

In some cases, applying these operators will not lead to any variations in the base-16 counts, which corresponds to the local symmetries of the first type. This occurs when the base-16 symbols, or operands, acted on by a given transposition operator share the same symbol in Alice's base-4 sequence. In all, there are $6$ operators that may each act on one of $4$ operands, leading to $24$ unique operations of this type. All base-16 sequences within the elementary state space $\varepsilon^a$ are related to one another through some composition of these 24 distinct transposition operations. An example of one such operation is as follows:

\begin{equation}
\scalebox{0.8}{$\begin{aligned}
\left(\begin{array}{c}
AA\\
BD\\
CA\\
CC\\
BB\\
AA
\end{array}\right)\xrightarrow{T^{AC}_{b2}}\left(\begin{array}{c}
AA\\
BD\\
CC\\
CA\\
BB\\
AA
\end{array}\right)
\end{aligned}$}
\label{eq: transpostion 1}
\end{equation}

Local symmetries of the second type are slightly more complicated. This is because any single transposition operator cannot simultaneously change the base-16 counts while also leaving the local quantum numbers invariant. An example of why this is the case can be seen in the following, where we pay special attention to the map relating Alice's and Bob's base-4 sequence: 

\begin{equation}
\scalebox{0.8}{$\begin{aligned}
\left(\begin{array}{c}
A\\
B\\
C\\
D\\
B\\
A
\end{array}\right)_{a1}\oplus\left(\begin{array}{c}
A\\
A\\
B\\
B\\
A\\
B
\end{array}\right)_{map}=\left(\begin{array}{c}
A\\
B\\
D\\
C\\
B\\
B
\end{array}\right)_{b2}\xrightarrow{T^{CD}_{b2}}
\left(\begin{array}{c}
A\\
B\\
C\\
D\\
B\\
A
\end{array}\right)_{a1}\oplus\left(\begin{array}{c}
A\\
A\\
A\\
A\\
A\\
B
\end{array}\right)_{map}=\left(\begin{array}{c}
A\\
B\\
C\\
D\\
B\\
B
\end{array}\right)_{b2}
\end{aligned}$}
\label{eq: transposition 2}
\end{equation}

Clearly, the base-4 counts associated with the map relating Alice's and Bob's base-4 sequence changed as a result of this operation. To avoid this, transposition operators associated with the second type of local symmetry must always act in pairs. For example, when the operator $T^{AB}_{b2}T^{CD}_{b2}$ acts on the base-16 symbols $AA$, $BB$, $CD$, and $DC$, we have the following:

\begin{equation}
\scalebox{0.8}{$\begin{aligned}
\left(\begin{array}{c}
A\\
B\\
C\\
D\\
B\\
A
\end{array}\right)_{a1}\oplus\left(\begin{array}{c}
A\\
A\\
B\\
B\\
A\\
B
\end{array}\right)_{map}=\left(\begin{array}{c}
A\\
B\\
D\\
C\\
B\\
B
\end{array}\right)_{b2}\xrightarrow[]{T^{AB}_{b2}T^{CD}_{b2}}
\left(\begin{array}{c}
A\\
B\\
C\\
D\\
B\\
A
\end{array}\right)_{a1}\oplus\left(\begin{array}{c}
B\\
B\\
A\\
A\\
A\\
B
\end{array}\right)_{map}=\left(\begin{array}{c}
B\\
A\\
C\\
D\\
B\\
B
\end{array}\right)_{b2}
\end{aligned}$}
\label{eq: transposition 3}
\end{equation}

As we claimed, the operation depicted in equation (\ref{eq: transposition 3}) indeed leads to the conservation of all local quantum numbers, while generating a change in the base-16 counts. In general, local symmetries of the second type can be expressed in terms of operators like $T^{AB}_{b2}T^{CD}_{b2}$, or ordered pairs of transpositions. In total, there are 36 unique pairs that can be created using the 6 basic operators. However, 6 of these are just repetitions of the same operator, which do not contribute to local symmetries of the second type:

\begin{equation}
    T^{AB}_{b2}T^{AB}_{b2}, \quad T^{AC}_{b2}T^{AC}_{b2}, \quad T^{AD}_{b2}T^{AD}_{b2}, \quad T^{BC}_{b2}T^{BC}_{b2}, \quad T^{BD}_{b2}T^{BD}_{b2}, \quad T^{CD}_{b2}T^{CD}_{b2}
\end{equation}

The remaining 30 can be divided into two classes, which are distinguished from one another through their behavior under commutation. The 6 pairs that commute are as follows:

\begin{equation}
    T^{AB}_{b2}T^{CD}_{b2}, \quad T^{CD}_{b2}T^{AB}_{b2}, \quad T^{AC}_{b2}T^{BD}_{b2}, \quad T^{BD}_{b2}T^{AC}_{b2}, \quad T^{AD}_{b2}T^{BC}_{b2}, \quad T^{BC}_{b2}T^{AD}_{b2}
    \label{eq: commuttative operators}
\end{equation}

Of course, due to their commutativity, these six pairs only represent 3 unique operators, bringing the previous total of 30 down to 27. The remaining 24 operators are those which do not commute, each of which involves the transposition of only 3 of the 4 base-4 symbols:

\begin{equation}
\begin{array}{cc}
     T^{BC}_{b2}T^{BD}_{b2}, \quad T^{BD}_{b2}T^{CD}_{b2}, \quad T^{CD}_{b2}T^{BC}_{b2}, \quad T^{BD}_{b2}T^{BC}_{b2}, \quad T^{CD}_{b2}T^{BD}_{b2}, \quad T^{BC}_{b2}T^{CD}_{b2}\\
     T^{AC}_{b2}T^{AD}_{b2}, \quad T^{AD}_{b2}T^{CD}_{b2}, \quad T^{CD}_{b2}T^{AC}_{b2}, \quad T^{AD}_{b2}T^{AC}_{b2}, \quad T^{CD}_{b2}T^{AD}_{b2}, \quad T^{AC}_{b2}T^{CD}_{b2}\\
     T^{AB}_{b2}T^{AD}_{b2}, \quad T^{AD}_{b2}T^{BD}_{b2}, \quad T^{BD}_{b2}T^{AB}_{b2}, \quad T^{AD}_{b2}T^{AB}_{b2}, \quad T^{BD}_{b2}T^{AD}_{b2}, \quad T^{AB}_{b2}T^{BD}_{b2}\\
     T^{AB}_{b2}T^{AC}_{b2}, \quad T^{AC}_{b2}T^{BC}_{b2}, \quad T^{BC}_{b2}T^{AB}_{b2}, \quad T^{AC}_{b2}T^{AB}_{b2}, \quad T^{BC}_{b2}T^{AC}_{b2}, \quad T^{AB}_{b2}T^{BC}_{b2}
\end{array}
\label{eq: non-commutative operators}
\end{equation}

The base-16 symbols being operated on, or the operands, are just as important as the operators themselves. To generate the desired local symmetry transformation, an operator must act on the appropriate collection of base-16 symbols. Our initial focus will be on the collections that obey a symmetry under the action of the commutative operators defined in equation (\ref{eq: commuttative operators}). Each of these collections consist of 4 base-16 symbols. For each one, we define a non-local quantum number consisting of a sum of the counts associated with each base-16 symbol. When the appropriate operator acts on the symbols associated with a given non-local quantum number, its numerical value will decrease by 1. Simultaneously, a different non-local quantum number will increase by 1. In general, non-local quantum numbers always come in pairs, such that the action of the appropriate operator leaves the sum of this pair invariant. These non-local quantum numbers are defined here:

\begin{equation}
\begin{array}{cc}
     \nu^{0}\equiv\frac{1}{4}( \widetilde{AA}+\widetilde{BB}+\widetilde{CD}+\widetilde{DC}), 
     \qquad \mu^{0}\equiv\frac{1}{4}( \widetilde{AB}+\widetilde{BA}+\widetilde{CC}+\widetilde{DD})\\
     \nu^{1}\equiv\frac{1}{4}( \widetilde{AD}+\widetilde{BC}+\widetilde{CA}+\widetilde{DB}), 
     \qquad \mu^{1}\equiv\frac{1}{4}( \widetilde{AC}+\widetilde{BD}+\widetilde{CB}+\widetilde{DA})\\
     \nu^{2}\equiv\frac{1}{4}( \widetilde{AA}+\widetilde{CC}+\widetilde{DB}+\widetilde{BD}),
     \qquad \mu^{2}\equiv\frac{1}{4}( \widetilde{AC}+\widetilde{CA}+\widetilde{DD}+\widetilde{BB})\\
     \nu^{3}\equiv\frac{1}{4}( \widetilde{BA}+\widetilde{DC}+\widetilde{AD}+\widetilde{CB}), 
     \qquad \mu^{3}\equiv\frac{1}{4}( \widetilde{BC}+\widetilde{DA}+\widetilde{AB}+\widetilde{CD})\\
     \nu^{4}\equiv\frac{1}{4}( \widetilde{AA}+\widetilde{DD}+\widetilde{BC}+\widetilde{CB}),
     \qquad \mu^{4}\equiv\frac{1}{4}( \widetilde{AD}+\widetilde{DA}+\widetilde{BB}+\widetilde{CC})\\
     \nu^{5}\equiv\frac{1}{4}( \widetilde{BA}+\widetilde{CD}+\widetilde{AC}+\widetilde{DB}), 
     \qquad \mu^{5}\equiv\frac{1}{4}( \widetilde{BD}+\widetilde{CA}+\widetilde{AB}+\widetilde{DC})
\end{array}
\label{eq: nu and mu}
\end{equation}

\begin{equation}
\begin{array}{cc}
     \kappa^{0}\equiv\frac{1}{4}( \widetilde{AA}+\widetilde{DC}+\widetilde{CB}+\widetilde{BD}), 
     \qquad \kappa^{1}\equiv\frac{1}{4}( \widetilde{AC}+\widetilde{DA}+\widetilde{CD}+\widetilde{BB})\\
     \kappa^{2}\equiv\frac{1}{4}( \widetilde{BA}+\widetilde{CC}+\widetilde{AD}+\widetilde{DB}), 
     \qquad \kappa^3\equiv\frac{1}{4}( \widetilde{BC}+\widetilde{CA}+\widetilde{AB}+\widetilde{DD})
\end{array}
\label{eq: kappa}
\end{equation}

\begin{equation}
\begin{array}{cc}
     \omega^0\equiv\frac{1}{4}( \widetilde{AA}+\widetilde{CD}+\widetilde{DB}+\widetilde{BC}), 
     \qquad \omega^1\equiv\frac{1}{4}( \widetilde{AD}+\widetilde{CA}+\widetilde{BB}+\widetilde{DC})\\
     \omega^2\equiv\frac{1}{4}( \widetilde{BA}+\widetilde{DD}+\widetilde{AC}+\widetilde{CB}), 
     \qquad \omega^3\equiv\frac{1}{4}( \widetilde{BD}+\widetilde{DA}+\widetilde{AB}+\widetilde{CC})\\
\end{array}
\label{eq: omega}
\end{equation}

These 20 non-local quantum numbers can be divided into two classes. The first, which are those defined in equation (\ref{eq: nu and mu}), are unique in that they are each associated with exactly 1 of the 3 commutative operators. These associations are as follows, where quantum numbers appear in parentheses along with their partner under the indicated operator:

\begin{equation}
\begin{array}{cc}
    T^{AB}_{b2}T^{CD}_{b2}: \qquad (\nu^0, \mu^0), \quad (\nu^1, \mu^1)\\
    T^{AC}_{b2}T^{BD}_{b2}: \qquad (\nu^2, \mu^2), \quad (\nu^3, \mu^3)\\
    T^{AD}_{b2}T^{BC}_{b2}: \qquad (\nu^4, \mu^4), \quad (\nu^5, \mu^5)
    \label{eq: nu mu operations}
\end{array}
\end{equation}

When one of these operations occurs, the only $\nu$'s or $\mu$'s impacted are those involved in the operation. Thus, these operations generate $\Delta\nu^i=\pm1$ and $\Delta\mu^i=\mp1$, while leaving all other $\nu$'s and $\mu$'s invariant. The second class of quantum numbers, which are those defined in equations (\ref{eq: kappa}) and (\ref{eq: omega}), are each associated with all 3 of the commutative operators. This implies that each of these non-local quantum numbers has three partners, one for each operator:

\begin{equation}
\begin{array}{cc}
    T^{AB}_{b2}T^{CD}_{b2}:\qquad(\kappa^0,\kappa^3),\quad (\kappa^1,\kappa^2),\quad(\omega^0, \omega^3), \quad (\omega^1,\omega^2)
    \\
    T^{AC}_{b2}T^{BD}_{b2}:\qquad(\kappa^0,\kappa^1),\quad (\kappa^2,\kappa^3),\quad(\omega^0, \omega^2), \quad (\omega^1,\omega^3)
    \\
    T^{AD}_{b2}T^{BC}_{b2}:\qquad(\kappa^0,\kappa^2),\quad (\kappa^1,\kappa^3),\quad(\omega^0, \omega^1), \quad (\omega^2,\omega^3)
    \\
\end{array}
\end{equation}

Like the operations in equation (\ref{eq: nu mu operations}), the only $\kappa$'s or $\omega$'s impacted by these operations are those actually involved, leaving all others invariant. This independence does not however span the two classes of quantum numbers we have outlined here. Any operation involving a $\nu$ and $\mu$ will induce a change of $\pm\frac{1}{2}$ in each of the $\kappa$'s and $\omega$'s. Similarly, operations involving $\kappa$'s or $\omega$'s will induce a change of $\pm\frac{1}{2}$ in 8 of the 12 $\nu$'s and $\mu$'s.

The final issue we must address in this section is that of the non-commutative operations. These operations each involve 3 base-16 symbols, rather than 4. The full set of non-local quantum numbers involved in these operations, of which there are 32, can be found in appendix \ref{sec: Non-commutative operands}. As with those previously discussed, each of these non-local quantum numbers has a partner under the action of a given non-commutative operator. In total, there are 16 partnerships of interest, each of which is associated with exactly six of the non-commutative operators in equation (\ref{eq: non-commutative operators}). As was the case with $\nu$'s/$\mu$'s and $\kappa$'s/$\omega$'s, the non-local quantum numbers defined in appendix \ref{sec: Non-commutative operands} lack independence from operations in which they are not directly involved. Likewise, the $\nu$'s/$\mu$'s and $\kappa$'s/$\omega$'s will vary under the action of the non-commutative operations. This complicated interdependence of the non-local quantum numbers introduced in this section will be relevant in the following section, where these quantum numbers will be used to drive interference. As we will see, care must be taken to ensure that the six non-local quantum numbers chosen to populate $\Lambda$ lead to a stable calculation. 

Needless to say, the local state spaces $L^a$ and $L^b$ have a rich set of symmetries. Though much of what has been discussed in this section remains physically obscure, several interesting features are of note. To begin with, we point out that a local state space is always associated with a specific choice of the local quantum numbers. In other words, one must assume a particular observable outcome for the experiment being modeled. Within this context, the symmetries of a local state space represent a means of organizing all of the unobservable states that contribute to that particular observable outcome. The fact that these symmetries come in two distinct forms is directly related to the nature of the hidden information. Symmetries of the first type arise because the information stored in the ordering of the symbols comprising a given sequence is hidden from observers. While symmetries of the second type arise because the numerical values of the non-local quantum numbers are hidden from observers. Put more simply, if no information were hidden from observers, no two ontic states would be indistinguishable within the proposed formalism. Thus, symmetries are directly connected to the concept of hidden information within this framework. 

\subsection{Interference \label{sec: Interference}}

In section \ref{sec: Indistinguishable states}, we identified quantum numbers that qualify as non-local within this formalism. We found them by studying the symmetries of local state spaces, which are associated with the transposition operators introduced in equations (\ref{eq: commuttative operators}-\ref{eq: non-commutative operators}). We may now use these results to elucidate the phenomenon of interference, or situations in which base-16 sequences within a local state space annihilate one another. 

Given two indistinguishable states from $L^a$ or $L^b$, which we label with the subscripts $0$ and $1$, the total difference between their non-local quantum numbers is as follows, where $\lambda$ is a 6-tuple:

\begin{equation}
    f(\lambda_0,\lambda_1)\equiv\sum_{k\in [0,5]}\lambda_0^{k}-\lambda_1^{k}
    \label{eq: interference term}
\end{equation}

 Within the local state spaces defined in equations (\ref{eq: local state space Alice}) and (\ref{eq: local state space Bob}), two states will annihilate one another if $f(\lambda_0,\lambda_1)$ is odd. These annihilations can be accounted for when calculating the post-interference cardinalities of the local state spaces $L^a$ and $L^b$ in the following way, where we use the notation $||^\circleddash$ to distinguish this operation from the pre-interference cardinality:

\begin{equation}
    |L^a(r,u,w)|^\circleddash=\left|\sum_{\lambda_i\in Q(r, u, w)}(-1)^{\sum_{k}(\lambda^k_0-\lambda^k_i)}|\varepsilon^a(r,u,w,\lambda_i)|\right|
    \label{eq: local cardinality Alice}
\end{equation}
\begin{equation}
    |L^b(r,u,w)|^\circleddash=\left|\sum_{\lambda_i\in Q(r, u, w)}(-1)^{\sum_{k}(\lambda^k_0-\lambda^k_i)}|\varepsilon^b(r,u,w,\lambda_i)|\right|
    \label{eq: local cardinality Bob}
\end{equation}

The absolute value operation on the right hand side of equations (\ref{eq: local cardinality Alice}) and (\ref{eq: local cardinality Bob}) is included to ensure that the resulting cardinalities are positive. This operation is necessary due to the arbitrary sign convention used in the interference term, which arises from the choice of $\lambda_0$. It should also be noted that the method of accounting for interference presented here differs from the one used in \cite{Powers2022, Powers2023}, though both methods are technically equivalent. 

The parity (odd/even) of the quantity defined in equation (\ref{eq: interference term}) contains information about the transposition operations through which two indistinguishable states are related. More precisely, it tells us the parity of the total number of these operations, which are those associated with the operators introduced in equations (\ref{eq: commuttative operators}) and (\ref{eq: non-commutative operators}). Exactly which transposition operations the parity of equation (\ref{eq: interference term}) is sensitive to depends solely on the 6 non-local quantum numbers chosen to populate $\Lambda$. Though a complete overview of these choices is beyond the scope of this work, we offer here a brief discussion of the three most relevant considerations when making this choice.

First, we must ensure that the variations in $\sum_{k}\lambda^k$ within a local state space are always integers. Otherwise, the cardinalities of local state spaces, as they are defined in equations (\ref{eq: local cardinality Alice}) and (\ref{eq: local cardinality Bob}), may become complex numbers. Non-local quantum numbers were defined in a way that ensures they always vary in integer increments under operations in which they are directly involved. However, the interdependence we discussed in section \ref{sec: Indistinguishable states} leads to the possibility of non-integer variations. These non-integer variations can be corrected for within the full sum by selecting combinations of non-local quantum numbers which collectively have integer variations. This is an important selection criteria when choosing non-local quantum numbers to populate $\Lambda$.

Second, the number and type of operations the parity of equation (\ref{eq: interference term}) is sensitive to must be considered. Again, this depends solely on the choice of the 6 non-local quantum numbers. Depending on this choice, operations involving the operators introduced in equations (\ref{eq: commuttative operators}) and (\ref{eq: non-commutative operators}) may or may not induce a variation in the sum $\sum_{k}\lambda^k$. Of the operations that do induce a variation in this quantity, some will lead to odd variations and some will lead to even variations. With respect to interference, it is the operations which yield odd variations in the sum $\sum_{k}\lambda^k$ which are of interest. Again, exactly which operations these will be depends on the choice of non-local quantum numbers, though no choice will ever encompass all possible operations.

Third, the choice of 6 non-local quantum numbers must form a complete set of base-16 quantum numbers when paired with the 10 local quantum numbers. This is a basic requirement of this formalism as it enables one to map quantum numbers to base-16 counts. The simplest way to verify that a given choice of 16 quantum numbers is complete is to express them in matrix form, and attempt to invert the matrix. If it is invertible, then this choice meets the condition of completeness. 

Even with these considerations in mind, there still remain many valid choices for the 6 non-local quantum numbers. For the models presented here, we make the following choice, which we store in the set $\Lambda$:

\begin{equation}
    \Lambda=( \nu^0,\nu^1,\nu^4,\nu^5, \kappa^1, \omega^2)
    \label{eq: non-local number choice}
\end{equation}

What has been presented in this section is by no means the complete picture associated with the phenomenon of interference. However, it is sufficient to justify our choice of the 6 non-local quantum numbers given in equation (\ref{eq: non-local number choice}). First, this choice guarantees that variations in the sum $\sum_{k}\lambda^k$ will always be integers. Second, this choice maximizes the sensitivity of equation (\ref{eq: interference term}) to the underlying transposition operations. This group of operations can be found in appendix \ref{sec: Chosen non-local operations}. Lastly, this choice will always form a complete set when paired with the 10 local quantum numbers. With this choice made, the only remaining task is to calculate the probability of observing a particular experimental outcome. 

\subsection{Calculating probabilities \label{sec: Calculating probabilities}}

We may now address the main point of this formalism, which is to calculate the probability $P(r|u)$. As discussed in section \ref{sec: Ontic state spaces}, the mathematical structure of interest is the joint state space $E^{a}(u)\Bar{\otimes}E^{b}(u)$. Recall that this joint state space includes all possible ordered pairs of states in $E^a(u)$ and $E^b(u)$ which share the same set of local quantum numbers. This implies that $E^{a}(u)\Bar{\otimes}E^{b}(u)$ may be expressed as a union over the Cartesian product of the local state spaces $L^a$ and $L^b$ like so:

\begin{equation}
    E^{a}(u)\Bar{\otimes}E^{b}(u)=\bigcup_{(r,w)\in Q(u)}\bigcup_{\chi_{a1}(w),\chi_{b2}(w)}L^a(r,u,w)\otimes L^b(r,u,w)
    \label{eq: joint state space}
\end{equation}

To calculate the probability of observing a particular choice of $r$, given $u$, we must first determine the post-interference cardinality of the Cartesian product space $L^a(r,u,w)\otimes L^b(r,u,w)$. By the combinatorial rule of product, this quantity is the scalar product of equations (\ref{eq: local cardinality Alice}) and (\ref{eq: local cardinality Bob}):

\begin{equation}
|L^a(r,u,w)\otimes L^b(r,u,w)|^\circleddash=|L^a(r,u,w)|^\circleddash|L^b(r,u,w)|^\circleddash
\label{eq: local product space cardinality}
\end{equation}

To account for the union over local state space configurations, the quantity in equation (\ref{eq: local product space cardinality}) must be scaled by the product of the following two expressions, where $\widetilde{\chi}_{a1}^i(r,u,w)$, for example, denotes the count associated with the $i^{th}$ base-4 symbol in the N-tuple $\chi_{a1}(w)$:

\begin{equation}
    G_{a1}(r,u,w) \equiv \frac{(\sum_{i}\widetilde{\chi}_{a1}^i(r,u,w))!}{\prod_{i}\widetilde{\chi}_{a1}^i(r,u,w)!}
    \label{eq: Alice local state space configurations}
\end{equation}
\begin{equation}
    G_{b2}(r,u,w) \equiv \frac{(\sum_{i}\widetilde{\chi}_{b2}^i(r,u,w))!}{\prod_{i}\widetilde{\chi}_{b2}^i(r,u,w)!}
    \label{eq: Bob local state space configurations}
\end{equation}

Thus, the total number of states associated with a particular choice of $r$, within the joint state space $E^{a}(u)\Bar{\otimes}E^{b}(u)$, can be found by summing over the product of equations (\ref{eq: local product space cardinality}-\ref{eq: Bob local state space configurations}) for all possible combinations of the local nuisance variables $w$:

\begin{equation}
    \Upsilon(r,u)=\sum_{w\in Q(r,u)}G_{a1}(r,u,w)G_{b2}(r,u,w)|L^a(r,u,w)|^\circleddash|L^b(r,u,w)|^\circleddash
    \label{eq: upsilon}
\end{equation}

This result is then normalized by the total post-interference cardinality of Alice's and Bob's joint state space, which can be obtained by summing over equation (\ref{eq: upsilon}) for all possible combinations of the random variables $r$:

\begin{equation}
    |E^{a}(u)\Bar{\otimes}E^{b}(u)|^\circleddash=\sum_{r\in Q(u)}\Upsilon(r,u)
\end{equation}

We are left with the following expression for the probability of observing a particular choice of $r$, given $u$:

\begin{equation}
    P(r|u)=\frac{\Upsilon(r,u)}{\sum_{r\in Q(u)}\Upsilon(r,u)}
    \label{eq: probability}
\end{equation}

We are now in a position to say something precise about the general nature of probabilities within the proposed formalism. First, they arise due to the presence of hidden information, implying that they are epistemic in origin. Though, it is important to note that the state spaces associated with distinct choices of observables are always disjoint. That is, no single base-16 sequence will ever appear in state spaces associated with two different sets of observables. This property distinguishes this proposal from the well known Spekkens toy model \cite{Spekkens2007}. It also implies that this approach is in good standing with the PBR no-go theorem \cite{Pusey2012}. 

Second, the act of measurement, or the performance of an experiment, may be interpreted as the revelation of the preexisting state of the system under study. Though, some modification to one's definition of a "system" may be necessary. By "system", we mean two causally connected events. Of course, this interpretation has many wide ranging implications, most of which are beyond the narrow scope of this paper. However, it is important to note that this particular version of realism does not compromise this framework's status with respect to quantum contextuality theorems like the Kochen-Specker theorem \cite{Kochen1968}. Rather, one may view the event centric picture as a natural implication of quantum contextuality and the closely related issue of Bell's inequality \cite{Clauser1969,Bell1964}. The latter point is made explicit in section \ref{sec: Bell test results}, where we report a clear violation of the equivalent CHSH inequality. 

The final point we make is related to the Frequentest interpretation of probabilities. The size of the conditioning variable $n$ has a significant impact on the size of state spaces, and thus on the continuity of probabilities. In the case that $n$ is small, it is conceivable that Alice and Bob could perform enough experiments as to have observed all possible ontic states, where we assume that no ontic state ever occurs more than once. In this case, the predicted probability and their measurement results should match exactly, which is in alignment with the Frequentest interpretation. A secondary consequence of a small $n$ is a loss of statistical independence, where we again assume that no state occurs more than once. That is, Alice's and Bob's past measurement outcomes would influence what they know about future experiments, ultimately leading to the possibility of completely deterministic experiments. In practice, however, we will assume that $n$ is extremely large, and possibly infinite. In the case that $n$ becomes infinite, the cardinality of ontic state spaces will generally become uncountable, with only a few extreme cases in which they remain countable. Because the measurement process is inherently discrete, and thus countable, statistical independence is ensured in this case, even if every state occurs only once.

\section{Models \label{sec: Models}}

The purpose of this section is to illustrate the manner in which the proposed framework and associated formalism can be used to model non-trivial physical systems, which are otherwise modeled using QM. In total, we will introduce three models involving spin, the last of which can be applied to a Bell test. Importantly, this section should not be viewed as a complete and general analysis of the modeled systems, nor the complete physical picture as it relates to the underlying formalism. Again, it is simply an illustration of how the proposed framework can be applied to specific physical systems. Further investigation and development is necessary before these models can be fully understood and validated.  

\subsection{Local quantum numbers \label{sec: Local quantum numbers}}

We begin the process of introducing physics models with a general discussion on local quantum numbers. These quantum numbers must be functions of the base-4 counts defined in equations (\ref{eq: base-4 to base-16 A}-\ref{eq: base-4 to base-16 Dmap}). Though there are many possible choices, certain compositions of base-4 counts are of special interest. Throughout this section, we will motivate and define the various local quantum numbers of interest in the models proposed here. 

The first local quantum number is the sequence length $n$. As discussed in section \ref{sec: Ontic state spaces}, this quantity may be interpreted as a continuity parameter, allowing probabilities to become arbitrarily smooth. This is because as $n$ grows larger, so too will the underlying state spaces. For example, the number of states in the set $S^4(n)$ is $16^n$. The sequence length $n$ may also be used to normalize other quantum numbers, allowing those to become arbitrarily smooth as well. Because all base-4 sequences within a given ontic state must be the same length, $n$ may be defined in terms of the base-4 counts associated with events, or maps:

\begin{equation}
\begin{array}{l}
    n \equiv \tilde{A}_{a1} + \tilde{B}_{a1} + \tilde{C}_{a1} + \tilde{D}_{a1}\\
    n \equiv \tilde{A}_{b2} + \tilde{B}_{b2} + \tilde{C}_{b2} + \tilde{D}_{b2} \\
    n \equiv \tilde{A}_{map} + \tilde{B}_{map} + \tilde{C}_{map} + \tilde{D}_{map}
\end{array}
\end{equation}

The next set of quantum numbers of interest, which we denote as $j_{a1}$ and $j_{b2}$, are proportional to the Hamming distance between the base-2 sequences comprising Alice's and Bob's event. Because the Hamming distance obeys a triangle inequality, we interpret these quantum numbers as the total spin observed at each event. These quantities are defined as follows, where the factor of $\frac{1}{2}$ is included to align with the traditional units of angular momentum:

\begin{equation}
j_{a1}\equiv\frac{\tilde{C}_{a1}+\tilde{D}_{a1}}{2}\,,\qquad j_{b2}\equiv\frac{\tilde{C}_{b2}+\tilde{D}_{b2}}{2}\label{eq: j}
\end{equation}

Closely related to $j_{a1}$ and $j_{b2}$ are the associated spin projection quantum numbers $m_{a1}$ and $m_{b2}$. Of course, these quantities must obey the ranges $-j_{a1}\leq m_{a1}\leq j_{a1}$ and $-j_{b2}\leq m_{b2}\leq j_{b2}$, respectively. This leaves us with two possible definitions, which are proportional to the differences $\tilde{C}-\tilde{D}$ or $\tilde{D}-\tilde{C}$. We choose the former, leading to the following definitions:

\begin{equation}
m_{a1}\equiv\frac{\tilde{C}_{a1}-\tilde{D}_{a1}}{2}\,,\qquad m_{b2}\equiv\frac{\tilde{C}_{b2}-\tilde{D}_{b2}}{2}\label{eq: m}
\end{equation}

The final quantum number associated with events does not yet have a clear physical interpretation. Because of this, the chosen definition should be viewed as operational and subject to change. We denote these provisional quantities as $l_{a1}$ and $l_{b2}$ and choose a definition which is analogous to $m_{a1}$ and $m_{b2}$:

\begin{equation}
l_{a1}\equiv\frac{\tilde{A}_{a1}-\tilde{B}_{a1}}{2}\,,\qquad l_{b2}\equiv\frac{\tilde{A}_{b2}-\tilde{B}_{b2}}{2}\label{eq: l}
\end{equation}

In all of the models considered here, the quantum numbers $n$, $j$, $m$, and $l$ will be those used to order the set $S^2(n)$, or the ontic state space for events. What remains is to define the quantum numbers associated with maps. For the time being, we denote these three quantum numbers using the generic symbols $\alpha_{map}$, $\beta_{map}$, and $\gamma_{map}$, which we define as follows:

\begin{equation}
    \alpha_{map} \equiv \tilde{B}_{map} + \tilde{C}_{map}
    \label{eq: alpha}
\end{equation}
\begin{equation}
    \beta_{map} \equiv \tilde{B}_{map} + \tilde{D}_{map}
    \label{eq: beta}
\end{equation}
\begin{equation}
    \gamma_{map} \equiv \tilde{C}_{map} + \tilde{D}_{map}
    \label{eq: gamma}
\end{equation}

When introducing each model, a physical interpretation will be offered for these abstract quantities when appropriate. There will also be several additional quantum numbers defined within those sections, which we refrain from introducing here for the sake of clarity. 

\subsection{Spin systems in rotated frames \label{sec: Spin systems in rotated frames}}

\begin{figure}
    \centering
    \includegraphics[scale=.75]{sketch_3.png}
    \caption{Two Stern-Gerlach detectors with a relative angle of rotation $\theta$ are assigned to separate observers named Alice and Bob. Within each detector occurs a single event, which share a direct causal connection to one another.}
    \label{fig: sketch 3}
\end{figure}

The physical system of interest in this section consists of two Stern-Gerlach detectors which may be rotated with respect to one another, as depicted in Figure \ref{fig: sketch 3}. In this case, an experiment consists of two events with a direct causal connection and a definite time ordering. That is, Alice's event precedes Bob's, or Bob's event precedes Alice's. For the purposes of this section, we will assume that Alice's event precedes Bob's. Within this context, one may imagine a single particle of spin $j$ first interacting with Alice's detector and then with Bob's. 

An obvious requirement of this model is that the spin quantum number associated with Alice's event be the same as the one associated with Bob's, or that $j_{a1}=j_{b2}$. This can be accomplished by choosing $\gamma_{map}=0$. In terms of the underlying base-16 counts, this constraint has the following impact, which can be inferred from equations (\ref{eq: base-4 to base-16 Cmap}), (\ref{eq: base-4 to base-16 Dmap}), and (\ref{eq: gamma}): 

\begin{equation}
    \widetilde{AC}+\widetilde{BD}+\widetilde{CA}+\widetilde{DB} + \widetilde{AD}+\widetilde{BC}+\widetilde{CB}+\widetilde{DA} = 0
\end{equation}

Thus, the proposed constraint reduces the total number of degrees of freedom within this model to 8, rather than 16. After applying this constraint, the quantum numbers associated with events can be expressed in terms of base-16 counts like so, where we note that $j_{a1}$ and $j_{b2}$ are now identical:

\begin{equation}
n=\widetilde{AA}+\widetilde{BB}+\widetilde{CC}+\widetilde{DD}+\widetilde{AB}+\widetilde{BA}+\widetilde{CD}+\widetilde{DC}
\label{eq: n rotation}
\end{equation}
\begin{equation}
j_{a1}=\frac{\widetilde{CC}+\widetilde{CD}+\widetilde{DC}+\widetilde{DD}}{2}\,,\qquad j_{b2}=\frac{\widetilde{CC}+\widetilde{DC}+\widetilde{CD}+\widetilde{DD}}{2}\label{eq: j rotation}
\end{equation}
\begin{equation}
m_{a1}=\frac{\widetilde{CC}+\widetilde{CD}-\widetilde{DC}-\widetilde{DD}}{2}\,,\qquad m_{b2}=\frac{\widetilde{CC}+\widetilde{DC}-\widetilde{CD}-\widetilde{DD}}{2}\label{eq: m rotation}
\end{equation}
\begin{equation}
l_{a1}=\frac{\widetilde{AA}+\widetilde{AB}-\widetilde{BA}-\widetilde{BB}}{2}\,,\qquad l_{b2}=\frac{\widetilde{AA}+\widetilde{BA}-\widetilde{AB}-\widetilde{BB}}{2}\label{eq: l rotation}
\end{equation}

A further consequence of setting $\gamma_{map}=0$ is that $\alpha_{map}=\beta_{map}=\tilde{B}_{map}$. This implies that the maps of interest within this model will only contain the symbols $A$ and $B$. An example of one such map is offered here: 

\begin{equation}
\scalebox{0.8}{$\begin{aligned}
\left(\begin{array}{c}
A\\
B\\
C\\
C\\
B\\
A
\end{array}\right)_{a1}\oplus\left(\begin{array}{c}
A\\
A\\
B\\
B\\
B\\
A
\end{array}\right)_{map}=\left(\begin{array}{c}
A\\
B\\
D\\
D\\
A\\
A
\end{array}\right)_{b2}\label{eq: rotation map}
\end{aligned}$}
\end{equation}

The relative angle of rotation between Alice's and Bob's detector is modeled here by the ratio $\frac{\tilde{B}_{map}}{n}$. The range of $\frac{\tilde{B}_{map}}{n}$ is $[0,1]$, which we map to $[0,\pi]$ by converting to radians. Note that this range can be made arbitrarily smooth by increasing $n$. We now offer the following definition of the angle $\theta$ within this model: 

\begin{equation}
    \theta \equiv \frac{\tilde{B}_{map}}{n} \pi
    \label{eq: theta}
\end{equation}

The 7 quantum numbers we have defined thus far constitute the full set of local quantum numbers necessary for the proposed model. The only remaining degree of freedom is the non-local quantum number $\nu^0$, or its partner $\mu^0$, as they are defined in equation (\ref{eq: nu and mu}). In fact, $\nu^0$ and $\mu^0$ are the only quantum numbers among those defined in equations (\ref{eq: nu and mu}-\ref{eq: omega}) that remain non-local under the constraint $\gamma_{map}=0$. In this case, we will use $\nu^0$ as the final quantum number. This quantity, which is closely related to $\mu_{a1,b2}$ in \cite{Powers2023}, is restated here:

\begin{equation}
    \nu^0 = \frac{1}{4}(\widetilde{AA} + \widetilde{BB} + \widetilde{CD} + \widetilde{DC})
\end{equation}

What remains is to decide which of the local quantum numbers to treat as random variables, conditioning variables, and local nuisance variables. Given the physical context, the random variable of interest is $m_{b2}$, or the spin projection quantum number associated with Bob's event. The conditioning variables are $n$, $j$, $m_{a1}$, and $\theta$, where $j=j_{a1}=j_{b2}$. Finally, the local nuisance variables will be $l_{a1}$ and $l_{b2}$. With these choices, we are left with the following definitions for the sets $R$, $U$, $W$, and $\Lambda$: 

\begin{equation}
    R = (m_{b2}), \;U = (n,j,m_{a1},\theta),\;W = (l_{a1},l_{b2}),\;\Lambda = (\nu^0)
\end{equation}


\subsection{Entangled spin systems \label{sec: Entangled spin systems}}

\begin{figure}
  \centering
  \includegraphics[scale=.75]{sketch_4.png}
  \caption{Two aligned Stern-Gerlach detectors are assigned to separate observers named Alice and Bob. Within each detector occurs a single event, which share an indirect causal connection to one another.}
  \label{fig: sketch 4}
\end{figure}

The physical system of interest in this section again consists of two Stern-Gerlach detectors, though we now require that these detectors be aligned ($\theta=0$), as depicted in Figure \ref{fig: sketch 4}. The experiments under consideration involve two events with an indirect causal connection and an indefinite time ordering. That is, Alice's event may precede Bob's, Bob's event may precede Alice's, or they may occur simultaneously. Of course, this implies that one can no longer imagine Alice's and Bob's events as being connected by a single particle, as in section \ref{sec: Spin systems in rotated frames}. Instead, they must be associated with particles that are related through a third event, which we may associate with an observer named Charlie.

As was the case in section \ref{sec: Spin systems in rotated frames}, the model introduced here involves setting one of the quantities defined in equations (\ref{eq: alpha}-\ref{eq: gamma}) equal to zero. For reasons that will become clear shortly, we set $\alpha_{map}=0$. In terms of the underlying base-16 counts, this constraint has the following impact, which can be inferred from equations (\ref{eq: base-4 to base-16 Bmap}), (\ref{eq: base-4 to base-16 Cmap}), and (\ref{eq: alpha}): 

\begin{equation}
    \widetilde{AB}+\widetilde{BA}+\widetilde{CD}+\widetilde{DC} + \widetilde{AC}+\widetilde{CA}+\widetilde{BD}+\widetilde{DB} = 0
\end{equation}

The proposed constraint reduces the total number of degrees of freedom within this model to 8, rather than 16. After applying this constraint, the quantum numbers associated with Alice's and Bob's events can be expressed in terms of base-16 counts like so:

\begin{equation}
n=\widetilde{AA}+\widetilde{BB}+\widetilde{CC}+\widetilde{DD}+\widetilde{AD}+\widetilde{BC}+\widetilde{CB}+\widetilde{DA}
\label{eq: n spin addition}
\end{equation}
\begin{equation}
j_{a1}=\frac{\widetilde{CC}+\widetilde{CB}+\widetilde{DA}+\widetilde{DD}}{2}\,,\qquad j_{b2}=\frac{\widetilde{CC}+\widetilde{BC}+\widetilde{AD}+\widetilde{DD}}{2}\label{eq: j spin addition}
\end{equation}
\begin{equation}
m_{a1}=\frac{\widetilde{CC}+\widetilde{CB}-\widetilde{DA}-\widetilde{DD}}{2}\,,\qquad 
m_{b2}=\frac{\widetilde{CC}+\widetilde{BC}-\widetilde{AD}-\widetilde{DD}}{2}\label{eq: m spin addition}
\end{equation}
\begin{equation}
l_{a1}=\frac{\widetilde{AA}+\widetilde{AD}-\widetilde{BC}-\widetilde{BB}}{2}\,,\qquad 
l_{b2}=\frac{\widetilde{AA}+\widetilde{DA}-\widetilde{CB}-\widetilde{BB}}{2}\label{eq: l spin addition}
\end{equation}

The constraint $\alpha_{map}=0$ also implies that $\beta_{map}=\gamma_{map}=\tilde{D}_{map}$. Thus, the maps of interest within this model will only contain the symbols $A$ and $D$. An example of one such map is offered here: 

\begin{equation}
\scalebox{0.8}{$\begin{aligned}
\left(\begin{array}{c}
A\\
B\\
C\\
C\\
B\\
A
\end{array}\right)_{a1}\oplus\left(\begin{array}{c}
A\\
A\\
D\\
A\\
A\\
D
\end{array}\right)_{map}=\left(\begin{array}{c}
A\\
B\\
B\\
C\\
B\\
D
\end{array}\right)_{b2}\label{eq: entanglement map}
\end{aligned}$}
\end{equation}

An important consequence of restricting maps to the symbols $A$ and $D$ is that the base-2 sequences carrying the subscripts $a$ and $b$ will always be identical. In other words, Alice's and Bob's events will always share a common reference sequence within this model. This leaves just three unique base-2 sequences in each ontic state within this model, which we may interpret as the vertices of a triangle:

\begin{equation}
    \begin{array}{cc}
         \includegraphics[scale=.8]{sketch_5.png}
    \end{array}
    \label{eq: triangle}
\end{equation}


Equations (\ref{eq: n spin addition}-\ref{eq: l spin addition}) represent the full set of local quantum numbers associated with two legs of this triangle. Those associated with the third leg are given here:


\begin{equation}
j_{12}\equiv\frac{\widetilde{BC}+\widetilde{DA}+\widetilde{AD}+\widetilde{CB}}{2}
\label{eq: j_12}
\end{equation}
\begin{equation}
m_{12}\equiv\frac{\widetilde{BC}+\widetilde{DA}-\widetilde{AD}-\widetilde{CB}}{2}
\label{eq: m_12}
\end{equation}
\begin{equation}
l_{12}\equiv\frac{\widetilde{AA}+\widetilde{CC}-\widetilde{BB}-\widetilde{DD}}{2}
\label{eq: l_12}
\end{equation}

The quantum numbers $j_{a1}$, $j_{b2}$, and $j_{12}$ are each proportional to the Hamming distance between the base-2 sequences indicated by their respective subscripts. This implies that we may think of them as the edge lengths for the triangle depicted in equation (\ref{eq: triangle}), which means that they obey the following selection rule of angular momentum addition:

\begin{equation}
    |j_{a1}-j_{b2}|\leq j_{12} \leq j_{a1}+j_{b2}
    \label{eq: triangle inequality}
\end{equation}

Of the three quantities defined in equations (\ref{eq: j_12}-\ref{eq: l_12}), $j_{12}$ is the only one that contains information not already included in equations (\ref{eq: j spin addition}-\ref{eq: l spin addition}). In fact, it is proportional to the number of $D$'s that appear in the map relating Alice's and Bob's events, which will be important in the following section. The quantum numbers $m_{12}$ and $l_{12}$, on the other hand, can each be expressed in terms of $m_{a1}$, $m_{b2}$, $l_{a1}$, and $l_{b2}$ in the following way:

\begin{equation}
m_{12}=-m_{a1}+m_{b2}=l_{b2}-l_{a1}\label{eq: m_12 identity}
\end{equation}
\begin{equation}
l_{12}=l_{a1}+m_{b2}=l_{b2}+m_{a1}\label{eq: l_12 identity}
\end{equation}

Equation (\ref{eq: m_12 identity}) contains the expression $m_{12}=-m_{a1}+m_{b2}$, which is reminiscent of conservation of angular momentum. In this form, the event at Bob's detector is associated with a composite system, while the event at Alice's detector is associated with a component of that system. To consider alternatives to this situation, we need a way to modify the signs of the spin projection quantum numbers. In general, this can be accomplished by commuting the base-2 sequences within Alice's and or Bob's events. Under these operations, the numerical values associated with the counts $\tilde{C}$ and $\tilde{D}$ are exchanged, while those associated with $\tilde{A}$ and $\tilde{B}$ are left invariant. By the definition of spin projection quantum numbers offered in equation (\ref{eq: m}), this exchange results in a sign change for $m$'s, while leaving all other quantum numbers associated with events unchanged. 

Unfortunately, commuting base-2 sequences within events complicates the quantum numbers associated with maps, as well as the non-local quantum numbers discussed in section \ref{sec: Indistinguishable states}. While this may indeed be the proper course in the long run, a simpler solution will suffice for the models of interest in this paper. In section \ref{sec: Local quantum numbers}, recall that we had two choices for the definition of the spin projection quantum numbers. Rather than choosing $\Tilde{C}-\tilde{D}$, we could have instead chosen $\Tilde{D}-\tilde{C}$. By allowing this choice to vary between the definition of Alice's and Bob's spin projection quantum numbers, we can capture the effect of the commutation operation, while leaving the quantum numbers we have already defined alone. For the purposes of this section and the next, we will choose this alternative definition for Alice's spin projection quantum number. We distinguish between these two definitions by swapping the subscripts associated with Alice's base-4 quantum numbers like so:

\begin{equation}
    j_{1a}=j_{a1}=\frac{\Tilde{D}_{a1}+\Tilde{C}_{a1}}{2}
\end{equation}
\begin{equation}
    m_{1a}=-m_{a1}=\frac{\Tilde{D}_{a1}-\Tilde{C}_{a1}}{2}
\end{equation}
\begin{equation}
    l_{1a}=l_{a1}=\frac{\Tilde{A}_{a1}+\Tilde{B}_{a1}}{2}
\end{equation}

With this modification, equations (\ref{eq: m_12}) and (\ref{eq: l_12}) become the following:

\begin{equation}
m_{12}=m_{1a}+m_{b2}=l_{b2}-l_{1a}\label{eq: m_12 identity-mod}
\end{equation}
\begin{equation}
l_{12}=l_{a1}+m_{b2}=l_{b2}-m_{1a}\label{eq: l_12 identity-mod}
\end{equation}

Thus, the spin addition selection rule $m_{12}=m_{1a}+m_{b2}$ now joins equation (\ref{eq: triangle inequality}) as an intrinsic property of the spin quantum numbers $j$ and $m$, under the proposed constraints. This result is central to the issue of entanglement, which demands a perfect correlation between Alice's and Bob's measurement outcomes when $\theta=0$. Importantly, equations (\ref{eq: m_12 identity-mod}) and (\ref{eq: l_12 identity-mod}) are both a direct consequence of Alice's and Bob's events sharing a common reference sequence. Note that this result can be traced all the way back to the choice to set $\alpha=0$.  

Before moving on, we again note that $j_{12}$ is actually proportional to the number of $D$'s that appear in the maps connecting Alice's and Bob's events. In this way, it is analogous to the quantum number $\theta$ introduced in the previous section. This association between the total spin of a composite system and the number of $D$'s that appear in maps will persist in the following section. In fact, it will supersede its association with the Hamming distance between the base-2 sequences carrying the subscripts $1$ and $2$. 

Due to the identities in equations (\ref{eq: m_12 identity-mod}) and (\ref{eq: l_12 identity-mod}), only seven of the ten quantum numbers introduced thus far are necessary to form a complete set. Our choices will be $n$, $j_{1a}$, $m_{1a}$, $j_{b2}$, $m_{b2}$, $j_{12}$, and $l_{12}$. The only remaining degree of freedom is the non-local quantum number $\nu^4$, or its partner $\mu^4$, as defined in equation (\ref{eq: nu and mu}). We will use $\nu^4$ in this case, which is closely related to $k$ in \cite{Powers2022}. This non-local quantum number is restated here:

\begin{equation}
    \nu^4 = \frac{1}{4}(\widetilde{AA}+\widetilde{DD}+\widetilde{BC}+\widetilde{CB})
\end{equation}

In this model, the random variables are $m_{1a}$ and $m_{b2}$, the conditioning variables are $n$, $j_{1a}$, $j_{b2}$, and $j_{12}$, and the local nuisance variable is $l_{12}$. Notably, the quantum number $m_{12}$ does not explicitly enter this probability calculation. However, it does constrain the allowed combinations of $m_{1a}$ and $m_{b2}$ through equation (\ref{eq: m_12 identity-mod}). We may interpret this as a consequence of $m_{12}$ being associated with a third event, which is not explicitly modeled. We now have the following definitions of the sets $R$, $U$, $W$, and $\Lambda$: 

\begin{equation}
    R = (m_{1a},m_{b2}), \;U = (n,j_{1a},j_{b2},j_{12}),\;W = (l_{12}),\; \Lambda = (\nu^4)
\end{equation}


\subsection{Bell Test \label{sec: Bell}}

\begin{figure}
    \centering
    \includegraphics[scale=.75]{sketch_6.png}
    \caption{Two Stern-Gerlach detectors with a relative angle of rotation $\theta_{ab}$ are assigned to separate observers named Alice and Bob. Within each detector occurs a single event, which share an indirect causal connection to one another.}
    \label{fig: sketch 6}
\end{figure}

In this section, we will combine the previous two models to produce one which will enable us to model a Bell test (Figure \ref{fig: sketch 6}). This can be accomplished by combining the maps used in the previous two sections under the addition mod2 operation. One of these maps is comprised of $A$'s and $D$'s and encodes the total spin of the composite system. The other is comprised of $A$'s  and $B$'s and encodes the relative rotation between the two detectors. These two maps sum to form a single unique base-4 sequence, which encodes both of these degrees of freedom. One example of this construction is given here:

\begin{equation}
\scalebox{0.8}{$\begin{aligned}
\left(\begin{array}{c}
A\\
B\\
B\\
A\\
B\\
A
\end{array}\right)\oplus\left(\begin{array}{c}
A\\
A\\
D\\
D\\
A\\
A
\end{array}\right)=\left(\begin{array}{c}
A\\
B\\
C\\
D\\
B\\
A
\end{array}\right)_{map}\label{eq: Bell rotation map}
\end{aligned}$}
\end{equation}

The number of $B$'s that appear in the map encoding rotation is equivalent to the sum $\tilde{B}+\tilde{C}=\alpha_{map}$ in the combined map. Similarly, the number of $D$'s in the map encoding the total spin of the composite system is equivalent to the sum $\tilde{D}+\tilde{C}=\gamma_{map}$ in the combined map. We are thus left with the following definitions of the angle between Alice's and Bob's detector and the total spin of the composite spin system within this model:


\begin{equation}
    j_{\gamma}=\frac{1}{2}\gamma_{map}
\end{equation}
\begin{equation}
    \theta_{ab}=\frac{\pi}{n}\alpha_{map}
    \label{eq: theta ab}
\end{equation}

The notation used for these more generally defined quantum numbers is intended to differentiate them from those defined in sections \ref{sec: Spin systems in rotated frames} and \ref{sec: Entangled spin systems}. The change from $j_{12}$ to $j_\gamma$ indicates that, once we allow for rotation, we can no longer interpret the total spin of the composite spin system as being associated with a Hamming distance. However, the triangle inequality defined in equation (\ref{eq: triangle inequality}) still holds. This is because the rotation maps, which involve $B$'s, always leave $j_{1a}$ and $j_{b2}$ invariant. 

On the other hand, the change from $\theta$ to $\theta_{ab}$ indicates that the new quantity is proportional to the Hamming distance between Alice's and Bob's reference sequences, as indicated by the subscripts. Actually, this is also the case for $\theta$, but in section \ref{sec: Spin systems in rotated frames}, $\theta$ was proportional to the Hamming distance between Alice's and Bob's complete base-4 sequences, rather than just the reference sequences.

These new quantum numbers can be paired with $n$, $j_{1a}$, $m_{1a}$, $l_{1a}$, $j_{b2}$, $m_{b2}$, and $l_{b2}$, which together constitute 9 of the 10 local quantum numbers associated with the experiments of interest in this section. The 10th local quantum number will be $\beta_{map}$, as defined in equation (\ref{eq: beta}). Like the quantum numbers $l_{1a}$ and $l_{b2}$, the physical interpretation of $\beta_{map}$ is unclear at this time. So, it will be treated as a nuisance variable within probability calculations, where its allowed range is as follows:

\begin{equation}
    |\frac{n}{\pi}\theta_{ab}-2j_{\gamma}|\leq \beta_{map} \leq \frac{n}{\pi}\theta_{ab}+2j_{\gamma}
\end{equation}

Unlike the models introduced in sections \ref{sec: Spin systems in rotated frames} and \ref{sec: Entangled spin systems}, this model requires all 6 non-local quantum numbers to form a complete set, which will be those chosen in equation (\ref{eq: non-local number choice}):

\begin{equation}
\begin{array}{cc}
    \nu^{0}\equiv\frac{1}{4}( \widetilde{AA}+\widetilde{BB}+\widetilde{CD}+\widetilde{DC})\\ 
    \nu^{1}\equiv\frac{1}{4}( \widetilde{AD}+\widetilde{BC}+\widetilde{CA}+\widetilde{DB})\\
    \nu^{4}\equiv\frac{1}{4}( \widetilde{AA}+\widetilde{DD}+\widetilde{BC}+\widetilde{CB})\\
    \nu^{5}\equiv\frac{1}{4}( \widetilde{BA}+\widetilde{CD}+\widetilde{AC}+\widetilde{DB})\\
    \kappa^{1}=\frac{1}{4}( \widetilde{AC}+\widetilde{DA}+\widetilde{CD}+\widetilde{BB})\\
    \omega^2=\frac{1}{4}( \widetilde{BA}+\widetilde{DD}+\widetilde{AC}+\widetilde{CB})
\end{array}
\end{equation}

The final issue we must address before modeling a Bell test concerns the spin projection quantum number for the composite system. Typically, Bell tests are performed on composite systems with total spin $0$. As a result, the spatial frame used to define this quantity is irrelevant in the resulting experiment. That is, only the relative angle between Alice's and Bob's detector has any physical significance. However, if the composite system has non-zero total spin, the spatial frame associated with this event does become relevant. In this case, we must keep track of Alice's and Bob's angle of rotation with respect to this third spatial frame, which we may associate with an observer named Charlie. Because Charlie's event is not explicitly included in this model, we are limited to scenarios in which this spatial orientation does not matter.  With this in mind, we may now address the definition of $m_\gamma$, which replaces $m_{12}$ as the spin projection quantum number for composite systems within this model. We begin by writing the definition of $m_\gamma$ under the constraint $\theta_{ab}=0$, which must be equivalent to $m_{12}$:

\begin{equation}
    (\theta_{ab}=0): m_\gamma=m_{12}=\frac{\widetilde{BC}+\widetilde{DA}-\widetilde{AD}-\widetilde{CB}}{2}
\end{equation}

In the case that $\theta_{ab}\neq0$, four additional base-16 counts will begin to contribute to $m_\gamma$. These are $\widetilde{AC}$, $\widetilde{CA}$, $\widetilde{BD}$, and $\widetilde{DB}$. To ensure that $-j_\gamma\leq m_\gamma \leq j_\gamma$, two of these additional base-16 counts must have a coefficient of $-\frac{1}{2}$ and two must have a coefficient of $+\frac{1}{2}$. This implies that we have six unique choices when defining $m_\gamma$, one for each combination of coefficients. Of these six choices, only two will yield the appropriate behavior for the full range of $\theta_{ab}$. These are defined as follows, where the distinguishing superscripts will be explained shortly:

\begin{equation}
    m^a_\gamma\equiv\frac{1}{2}(\widetilde{BC}+\widetilde{DA}+\widetilde{BD}+\widetilde{DB}-\widetilde{AD}-\widetilde{CB}-\widetilde{AC}-\widetilde{CA})
    \label{eq: m_gamma alice}
\end{equation}
\begin{equation}
    m^b_\gamma\equiv\frac{1}{2}(\widetilde{BC}+\widetilde{DA}+\widetilde{AC}+\widetilde{CA}-\widetilde{AD}-\widetilde{CB}-\widetilde{BD}-\widetilde{DB})
    \label{eq: m_gamma bob}
\end{equation}

The difference between these two quantities can be best understood by considering cases in which they are non-zero. As an example, consider the case in which $j_\gamma=1$, $m^a_\gamma=m^b_\gamma=+1$, $j_{1a}=j_{b2}=+\frac{1}{2}$. When $\theta_{ab}=0$, the only non-zero base-16 counts involved in the definition of $m^a_\gamma$ and $m^b_\gamma$ are $\widetilde{BC}$ and $\widetilde{DA}$, which both equal 1. When $\theta_{ab}=\pi$, these counts go to zero. To ensure that $m^a_\gamma$ and $m^b_\gamma$ are preserved under this limit, we must have $\widetilde{BD}=\widetilde{DB}=1$ ($m^a_\gamma=+1$) or $\widetilde{AC}=\widetilde{CA}=1$ ($m^b_\gamma=+1$). Therefore, when $\theta_{ab}=\pi$, $m^a_\gamma=+1$ implies $m_{1a}=+\frac{1}{2}$ and $m_{b2}=-\frac{1}{2}$ with certainty, while $m^b_\gamma=+1$ implies $m_{1a}=-\frac{1}{2}$ and $m_{b2}=+\frac{1}{2}$ with certainty.

The distinction between $m^a_\gamma$ and $m^b_\gamma$ outlined above actually holds for arbitrary rotations. If $m^a_\gamma$ is used, the spin projection quantum number observed at Alice's detector will always be $+\frac{1}{2}$, while Bob's may vary. Alternatively, using $m^b_\gamma$ will ensure $m_{b2}=+\frac{1}{2}$, while $m_{1a}$ may vary. Physically, the distinction between $m^a_\gamma$ and $m^b_\gamma$ can be understood in terms of spatial frames. That is, using $m^a_\gamma$ implies the spin projection quantum number for the composite system has been defined with respect to Alice's spatial frame, while $m^b_\gamma$ implies it has been defined with respect to Bob's. Of course, when $m^a_\gamma=m^b_\gamma=0$, this distinction becomes unimportant. 

As in section \ref{sec: Entangled spin systems}, the spin projection quantum number for the composite system will not explicitly enter the probability calculation. Instead, it will act as a constraint on the quantum numbers that do enter this calculation. In terms of these quantum numbers, $m^a_\gamma$ and $m^b_\gamma$ can be expressed as follows:

\begin{equation}
    m^a_\gamma = \frac{n}{4}+\frac{m_{1a}}{2}-\frac{l_{1a}}{2}-2(2\omega^2+\nu^0+\nu^1-\nu^4-\nu^5)
    \label{eq: m_gamma alice alt}
\end{equation}
\begin{equation}
    m^b_\gamma = -\frac{n}{4}+\frac{m_{b2}}{2}+\frac{l_{b2}}{2}+\frac{1}{2}(\frac{n}{\pi}\theta_{ab}-2j_\gamma)+2(2\kappa^1-\nu^0+\nu^1+\nu^4-\nu^5)
    \label{eq: m_gamma bob alt}
\end{equation}

In the following section, we will use $m^a_\gamma$ to constrain the permitted ranges of the relevant quantum numbers. Thus, equation (\ref{eq: m_gamma alice alt}) will supersede equation (\ref{eq: m_12 identity-mod}) as the selection rule governing this more general calculation. The only remaining task of this section is the assignment of quantum numbers to their respective categories, which are as follows:

\begin{equation}
    R = (m_{1a},m_{b2}), \;U = (n,j_{1a},j_{b2},\theta_{ab},j_{\gamma}),\;W = (l_{1a},l_{b2},\beta_{map}),\;\Lambda = ( \nu^0,\nu^1,\nu^4,\nu^5, \kappa^1, \omega^2)
\end{equation}


\section{Results \label{sec: results}}


In this section, we compare the output of equation (\ref{eq: probability}) to the predictions of QM  for each of the models introduced in section \ref{sec: Models}. In general, agreement with QM improves as $n$ increases. Though, due to computational limitations, the maximum value of $n$ considered here is $100$.  For the model introduced in section \ref{sec: Spin systems in rotated frames}, we consider systems with total spin $\frac{1}{2}$ and $1$. We find a meaningful deviation from QM which does not tend to zero for large $n$. Instead, the difference between the two models appears to be systematic in nature. We correct for this by modifying the definition of the angle between Alice's and Bob's spatial frames, which requires the introduction of a tuning parameter ($x$). For the model introduced in section \ref{sec: Entangled spin systems}, we consider constituent spin systems with total spin $\frac{1}{2}$ and $1$, and composite spin systems with total spin $0$, $1$, and $2$. In all but one case, exact agreement is found between QM and the proposed model. In the remaining case, the difference between the two models tends to zero as $n$ becomes large. Finally, the model introduced in section \ref{sec: Bell} is used to perform a CHSH calculation. For all physical systems considered here, with $n=100$ and $x=0.1377$, deviations between the proposed models and QM are less than $1\%$. 

\subsection{Results: spin systems in rotated frames \label{sec: rotated frames results}}

\begin{figure}
  \begin{centering}
      \subfigure{\includegraphics[scale=1]{Figure_1.pdf}}
      \subfigure{\includegraphics[scale=1]{Figure_2.pdf}}
  \end{centering}
  \caption{A comparison of models for $\theta$ (Left) and $\theta'$ (right), where $n=100$, $j=\frac{1}{2}$, $m_{a1}=+\frac{1}{2}$, and $|\Delta|$ is the magnitude of the difference between equations (\ref{eq: Wigner}) and (\ref{eq: Probability for rotation}).\label{fig: j=1/2 both thetas}}
\end{figure}

\begin{figure}
  \begin{centering}
      \subfigure{\includegraphics[scale=1]{Figure_3.pdf}}
      \subfigure{\includegraphics[scale=1]{Figure_4.pdf}}
  \end{centering}
  \caption{A comparison of models for $\theta$ (Left) and $\theta'$ (right), where $n=100$, $j=1$, $m_{a1}=+1$, and $|\Delta|$ is the magnitude of the difference between equations (\ref{eq: Wigner}) and (\ref{eq: Probability for rotation}).\label{fig: j=1 ma1=+1 both thetas}}
\end{figure}

\begin{figure}
  \begin{centering}
      \subfigure{\includegraphics[scale=1]{Figure_5.pdf}}
      \subfigure{\includegraphics[scale=1]{Figure_6.pdf}}
  \end{centering}
  \caption{A comparison of models for $\theta$ (Left) and $\theta'$ (right), where $n=100$, $j=1$, $m_{a1}=0$, and $|\Delta|$ is the magnitude of the difference between equations (\ref{eq: Wigner}) and (\ref{eq: Probability for rotation}).\label{fig: j=1 ma1=0 both thetas}}
\end{figure}

The physical system of interest in this section is a particle with total spin $j$, which first interacts with Alice's Stern-Gerlach detector, followed by Bob's. The probability we are interested in calculating is for the spin projection quantum number observed by Bob, which we denote as $m_{b2}$. In QM, this probability is given by the square of Wigner's d-matrix formula, which is a function of $j$, $m_{a1}$, $m_{b2}$, and $\theta$:

\begin{equation}
    \scalebox{1}{$\begin{aligned}
d_{m_{b2},m_{a1}}^{j}(\theta)=\sum_{q}(-1)^{m_{b2}-m_{a1}+q}\frac{(j+m_{a1})!(j-m_{a1})!(j+m_{b2})!(j-m_{b2})!}{(j+m_{a1}-q)!q!(m_{b2}-m_{a1}+q)!(j-m_{b2}-q)!}\\
\times\left(cos(\frac{\theta}{2})\right)^{2j+m_{a1}-m_{b2}-2q}\left(sin(\frac{\theta}{2})\right)^{m_{b2}-m_{a1}+2q}
\end{aligned}$}
\label{eq: Wigner}
\end{equation}

For the model proposed in section \ref{sec: Spin systems in rotated frames}, the expression which generates the probability of observing a particular value of $m_{b2}$ is given by the following, where $l_{a1}$, $l_{b2}$, and $\nu^0$ are summed over as outlined in section \ref{sec: Calculating probabilities}:

\begin{equation}
    P(m_{b2}|n,j,m_{a1},\theta)=\frac{\Upsilon(n,j,m_{a1},m_{b2},\theta)}{\sum_{m_{b2}\in Q(n,j,m_{a1},\theta)}\Upsilon(n,j,m_{a1},m_{b2},\theta)}
    \label{eq: Probability for rotation}
\end{equation}

In Figure \ref{fig: j=1/2 both thetas}(Left), a comparison of these two expressions is offered for $j=\frac{1}{2}$, $n=100$, and a full range of $\theta$. While there is a clear functional similarity between these two expressions, the curve associated with equation (\ref{eq: Probability for rotation}) appears to have a systematic bias towards $\frac{\pi}{2}$, which is maximal at $\frac{\pi}{4}$ and $\frac{3\pi}{4}$. Motivated by this observation, we propose the following alternative definition of $\theta$, as well as a proposed value for the tuning parameter $x$:

\begin{equation}
    \theta' \equiv \frac{\tilde{B}_{map}}{n} \pi - x\sin{\frac{2\tilde{B}_{map}}{n} \pi}
    \label{eq: theta_prime}
\end{equation}
\begin{equation}
    x=0.1377
    \label{eq: x value}
\end{equation}

In Figure \ref{fig: j=1/2 both thetas}(Right), a comparison between QM and the proposed model is offered for $\theta'$. In this case, deviations between QM and the proposed model are on the order of $0.1\%$. In Figures \ref{fig: j=1 ma1=+1 both thetas} and \ref{fig: j=1 ma1=0 both thetas}, identical comparisons are offered for physical systems with total spin $1$ and spin projections $m_{a1}=+1$ and $m_{a1}=0$, respectively. We again see in these cases that $\theta'$ yields much better agreement between the proposed model and QM. Though, the degree to which they differ is now on the order of $1\%$, rather than $0.1\%$, as in the spin $\frac{1}{2}$ case. Generally, for fixed $n$, the deviation between the proposed model and QM will increase as a function of total spin. A sample calculation for this model can be found in appendix \ref{sec: sample rotated}.

\subsection{Results: entangled spin systems  \label{sec: entangled results}}

The physical system of interest in this section involves two spin systems with total spin $j_{1a}$ and $j_{b2}$, which are members of a composite system with total spin $j_{12}$. Each constituent spin system interacts with either Alice's or Bob's Stern-Gerlach detector, which are spatially aligned. The probability we are interested in calculating is for the spin projection quantum numbers observed by Alice and Bob, which we denote as $m_{1a}$ and $m_{b2}$. In QM, this probability is given by the square of the Clebsch-Gordan coefficient, which is a function of $j_{1a}$, $m_{1a}$, $j_{b2}$, $m_{b2}$, $j_{12}$, and $m_{12}$:

\begin{equation}
\begin{array}{cc}
  <j_{1a}j_{b2}j_{12}m_{12}|j_{1a}j_{b2}m_{1a}m_{b2}>=\sqrt{\frac{(2j_{12}+1)(j_{1a}+j_{b2}-j_{12})!(j_{12}+j_{1a}-j_{b2})!(j_{12}+j_{b2}-j_{1a})!}{(j_{1a}+j_{b2}+j_{12}+1)!}}\\
\cdot\sum_{z}(-1)^{z}\frac{\sqrt{(j_{1a}+m_{1a})!(j_{1a}-m_{1a})!(j_{b2}+m_{b2})!(j_{b2}-m_{b2})!(j_{12}+m_{12})!(j_{12}-m_{12})!}}{z!(j_{1a}+j_{b2}-j_{12}-z)!(j_{1a}-m_{1a}-z)!(j_{b2}+m_{b2}-z)!(j_{12}-j_{b2}+m_{1a}+z)!(j_{12}-j_{1a}-m_{b2}+z)!}
\end{array}
  \label{eq: Original CGs}
\end{equation}

For the model proposed in section \ref{sec: Entangled spin systems}, the expression which generates the probability of observing a particular combination of $m_{1a}$ and $m_{b2}$ is given by the following, where $l_{12}$ and $\nu^4$ are summed over as outlined in section \ref{sec: Calculating probabilities}, and $m_{12}$ acts as a constraint:

\begin{equation}
    P(m_{1a},m_{b2}|n,j_{1a},j_{b2},j_{12})=\frac{\Upsilon(n,j_{1a},m_{1a},j_{b2},m_{b2},j_{12})}{\sum_{m_{1a},m_{b2}\in Q(n,j_{1a},j_{b2},j_{12})}\Upsilon(n,j_{1a},m_{1a},j_{b2},m_{b2},j_{12})}
    \label{eq: Probability for entangled}
\end{equation}

In the case that $j_{1a}=j_{b2}=\frac{1}{2}$ and $n\geq 2$, the proposed model and QM agree exactly for $j_{12}=1$ and $j_{12}=0$. For $j_{1a}=j_{b2}=1$ and $n\geq 4$, perfect agreement between the two models is again found for $j_{12}=2$, $j_{12}=1$, and $j_{12}=0$, in all but one case. The only case in which a difference exists is for $j_{12}=1$ and $m_{12}=0$. However, this difference tends to zero as $n$ becomes large, as depicted in Figure \ref{fig: CGC compare}. Thus, with $n=100$, the differences between the proposed model and QM are much less than $1\%$ for the systems considered here. A sample calculation for this model can be found in appendix \ref{sec: sample entangled}.

\begin{figure}
    \centering
    \includegraphics{Figure_7.pdf}
    \caption{A comparison of equations (\ref{eq: Original CGs}) and (\ref{eq: Probability for entangled}) over the range $4 \leq n \leq 100$, where $j_{1a}=1$, $j_{b2}=1$, $j_{12}=1$, $m_{12}=0$, and $|\Delta|$ is the magnitude of the difference between the two models. \label{fig: CGC compare}}
\end{figure}

\subsection{Results: Bell test \label{sec: Bell test results}}

The physical system of interest in this section involves two spin systems with total spin $j_{1a}=j_{b2}=\frac{1}{2}$, which are members of a composite system with total spin $j_{\gamma}=0$. Each constituent spin system interacts with either Alice's or Bob's Stern-Gerlach detector, which may be rotated with respect to one another. The model introduced in section \ref{sec: Bell} can be applied to this physical system to calculate the expectation values of the product $4m_{1a}m_{b2}$, as a function of the angle between Alice's and Bob's detectors. These expectation values can then be used to check if the proposed model indeed violates the CHSH inequality \cite{Clauser1969}, and if so, how it compares to the prediction of QM. We begin this calculation by writing the general expressions for the probability that Alice and Bob observe a particular combination of $m_{1a}$ and $m_{b2}$, where $l_{1a}$, $l_{b2}$, $\beta_{map}$, $\nu^0$, $\nu^1$, $\nu^4$, $\nu^5$, $\kappa^1$, and $\omega^2$ are summed over as outlined in section \ref{sec: Calculating probabilities}, and $m^a_{\gamma}$ acts as a constraint:

\begin{equation}
    P(m_{1a},m_{b2}|n,j_{1a},j_{b2},\theta_{ab},j_{\gamma})=\frac{\Upsilon(n,j_{1a},m_{1a},j_{b2},m_{b2},\theta_{ab},j_{\gamma})}{\sum_{m_{1a},m_{b2}\in Q(n,j_{1a},j_{b2},\theta_{ab},j_{\gamma})}\Upsilon(n,j_{1a},m_{1a},j_{b2},m_{b2},\theta_{ab},j_{\gamma})}
    \label{eq: Probability for Bell}
\end{equation}

Given $j_{1a}=j_{b2}=\frac{1}{2}$, $j_\gamma=m^a_\gamma=0$, and some choice of the angle $\theta_{ab}$, the expectation value of the product $4m_{1a}m_{b2}$ is given by the following:

\begin{equation}
    \mathcal{E}(n,\theta_{ab})=\sum_{m_{1a},m_{b2}}4m_{1a}m_{b2}P(m_{1a},m_{b2}|n,\frac{1}{2},\frac{1}{2},\theta_{ab},0)
\end{equation}

This expression can then be used to calculate the CHSH test statistic like so, where each angle is associated with a unique combination of detector orientations:

\begin{equation}
    \mathcal{S}(n,\theta_{ab},\theta_{ab'},\theta_{a'b},\theta_{a'b'})=|\mathcal{E}(n,\theta_{ab})-\mathcal{E}(n,\theta_{ab'})+\mathcal{E}(n,\theta_{a'b})+\mathcal{E}(n,\theta_{a'b'})|
    \label{eq: test statistic}
\end{equation}

In a Bell test, Alice and Bob each have two possible detector orientations, which are typically defined with respect to some fixed coordinate system. However, it is only the relative angle between the two detectors that is physically meaningful when $m^a_\gamma=0$, as discussed at the end of section \ref{sec: Bell}. This relative angle is precisely the quantity modeled by $\theta_{ab}$, as defined in equation (\ref{eq: theta ab}). The detector orientations we shall use in this analysis, as defined with respect to some fixed coordinate system, are as follows:

\begin{equation}
    \theta_a=0^{\circ}, \; \theta_{a'}=90^{\circ}, \; \theta_b=45^{\circ}, \; \theta_{b'}=135^{\circ}
\end{equation}

With these detector orientations, the four possible angles between Alice's and Bob's detectors are: 

\begin{equation}
    \theta_{ab}=45^{\circ}, \; \theta_{ab'}=135^{\circ}, \; \theta_{a'b}=45^{\circ}, \; \theta_{a'b'}=45^{\circ}
    \label{eq: angles}
\end{equation}

For this selection of angles, and $n=100$, the expression in equation (\ref{eq: test statistic}) yields the following:

\begin{equation}
     \mathcal{S}(100,45^{\circ},135^{\circ},45^{\circ},45^{\circ})=3.190
     \label{eq: x=0 CHSH}
\end{equation}

\begin{figure}
  \begin{centering}
      \subfigure{\includegraphics[scale=1]{Figure_8.pdf}}
      \subfigure{\includegraphics[scale=1]{Figure_9.pdf}}
  \end{centering}
  \caption{A plot of equation (\ref{eq: test statistic}) for $x=0$ (Left) and $x=0.1377$ (Right) for every value of $n$ in the range $1 \leq n \leq 100$ and the angles specified in equation (\ref{eq: angles}). In both plots, the red line indicates the Tsirelson bound of $2\sqrt{2}$.\label{fig: CHSH as a function of n both thetas}}
\end{figure}

Thus, the proposed model indeed violates the CHSH inequality ($\mathcal{S}\leq 2$). However, it also violates the Tsirelson bound of $2\sqrt{2}$, or the maximum value predicted by QM. The difference between the predictions of QM and the proposed model can be traced back to the difference seen in Figure \ref{fig: j=1/2 both thetas}(Left). Motivated by the results of section \ref{sec: rotated frames results}, we propose the following alternative definition of $\theta_{ab}$:

\begin{equation}
    \theta_{ab}'\equiv \frac{\alpha_{map}}{n}\pi - x\sin{\frac{\alpha_{map}}{n}2\pi}
    \label{eq: theta_ab prime}
\end{equation}

In the case that $x=0.1377$, as in section \ref{sec: rotated frames results}, equation (\ref{eq: x=0 CHSH}) yields the following:

\begin{equation}
    \mathcal{S}'(100,45^{\circ},135^{\circ},45^{\circ},45^{\circ})=2.846
     \label{eq: x=0.1377 CHSH}
\end{equation}

Due to the granularity of the angles $\theta_{ab}$ and $\theta_{ab}'$, the precise behavior of the CHSH test statistic as $n$ becomes large can be difficult to identify by only studying a small range of $n$'s. In Figures \ref{fig: CHSH as a function of n both thetas}(Left) and \ref{fig: CHSH as a function of n both thetas}(Right), the CHSH test statistic is plotted for the range $1\leq n \leq 100$ with $x=0$ (Left) and $x=0.1377$ (right). In both cases, the CHSH test statistic tends towards a constant value. For $x=0.1377$, this constant value is in the neighborhood of $2\sqrt{2}$, or the maximum value predicted by QM. A sample calculation for this model can be found in appendix \ref{sec: sample bell}.

\section{Discussion \label{sec: Discussion}}

We have presented here a fresh approach to modeling physics, which assumes the existence of a causal network consisting only of events and their relationships. The focus of this work has been only the simplest possible elements of this network, which each consist of two causally connected events. To model these exceedingly simple physical elements, we chose equally simple mathematical elements. Using only ordered lists of $0$'s and $1$'s, a robust formalism was developed and shown to have many non-trivial features of interest with respect to modeling physics. These include non-determinism, a rich and diverse set of symmetries, interference, conservation laws, and much more. Perhaps equally important, this formalism is accompanied by a vivid conceptual picture, consisting of causal networks of events, observers, their statistical ensembles, and the evolution of those ensembles. 

In section \ref{sec:Introduction}, we began by advocating for the event centric picture of nature. Though this approach has many merits, it was not assumed at the onset of this research program. The only assumption made was that $0$'s and $1$'s must be the simplest possible modeling tool a physicist could employ. By carefully studying these simple building blocks, we gradually identified features which appeared to be related to physics. Only after a long series of incremental advancements towards a robust physics model \cite{Powers2022, Powers2023}, did the presented framework become clear. Thus, the event centric picture was not chosen or assumed, it simply emerged. Though one should certainly not mistake this as evidence of scientific validity, the mathematical and conceptual simplicity of this picture is noteworthy. 

While the results presented here certainly justify further inquiry, there remains a long road ahead. Ultimately, the event centric approach must reproduce the modeling capabilities of quantum field theory and general relativity, simultaneously. While this remains the long term goal, we believe that the work presented here may already be of interest to those in the fields of quantum foundations, information theory, and even quantum computing. After all, the systems we have modeled here are the subject of intense study within these fields. Not only that, but the formalism which supports these models is built from classical bits. This may help clarify the relationship between classical and quantum information theory, for example.  

We now look ahead to the next few steps in the development of this framework. First and foremost, we must improve our understanding of how it fits within the broader literature on event centric approaches. Undoubtedly, the most efficient path to a robust event centric physics model will involve the confluence of many independent ideas. With that being said, there is no shortage of open questions and opportunities for further development within the framework presented here. For instance, a more complete treatment of statistical ensembles and their evolution may yield a connection to the wave equations employed in existing models. Also of interest for future work is a more in depth study of the symmetries of local state spaces, as this may clarify how the gauge symmetries in quantum field theory arise within this framework. Lastly, models for network elements containing more than two events is of significant interest. Of course, this list is far from complete. Though the proposed formalism is exceedingly simple, one can see that it still affords a great deal of complexity.
 
The mathematical simplicity, conceptual clarity, and close agreement with the predictions of established theory make the framework presented here of interest for further inquiry. Though many non-trivial advancements separate the present state of this research program from its ultimate goal, the achievements of other event centric approaches indicate that such an outcome may well be possible in the long run. In the near term, we are confident that further development will continue to shed light on the foundations of quantum theory, with potential implications for both pure and applied research programs. We look forward to these developments and to productive collaboration with the physics community at large. 


\section*{Acknowledgments}
We would like to thank Djordje Minic, Tatsu Takeuchi, Lauren Hay, and Omar Elsherif for comments on previous versions of this work as well as many helpful discussions. We also thank Emily Powers for figure
designs. D.S. is partially supported by the US National Science Foundation, under Grant no. PHY-2014021. 

\bibliography{references}

\appendix

\section{Non-commutative operands \label{sec: Non-commutative operands}}

The non-local quantum numbers associated with the non-commutative operators given in equation (\ref{eq: non-commutative operators}) are as follows: 

\begin{equation}
\begin{array}{ll}
     \rho^{0}\equiv\frac{1}{3}( \widetilde{BD}+\widetilde{CB}+\widetilde{DC}) & \eta^{0}\equiv\frac{1}{3}( \widetilde{BC}+\widetilde{CD}+\widetilde{DB})\\
     \rho^{1}\equiv\frac{1}{3}( \widetilde{AA}+\widetilde{BD}+\widetilde{CB}) & \eta^{1}\equiv\frac{1}{3}( \widetilde{AD}+\widetilde{BB}+\widetilde{CA})\\
     \rho^{2}\equiv\frac{1}{3}( \widetilde{AA}+\widetilde{BD}+\widetilde{DC}) & \eta^{2}\equiv\frac{1}{3}( \widetilde{AC}+\widetilde{BA}+\widetilde{DD})\\
     \rho^{3}\equiv\frac{1}{3}( \widetilde{AA}+\widetilde{CB}+\widetilde{DC}) & \eta^{3}\equiv\frac{1}{3}( \widetilde{AB}+\widetilde{CC}+\widetilde{DA})\\
     
     \rho^{4}\equiv\frac{1}{3}( \widetilde{AC}+\widetilde{BB}+\widetilde{DA}) & \eta^{4}\equiv\frac{1}{3}( \widetilde{AA}+\widetilde{BC}+\widetilde{DB})\\
     \rho^{5}\equiv\frac{1}{3}( \widetilde{AC}+\widetilde{CD}+\widetilde{DA}) & \eta^{5}\equiv\frac{1}{3}( \widetilde{AD}+\widetilde{CA}+\widetilde{DC})\\
     \rho^{6}\equiv\frac{1}{3}( \widetilde{BB}+\widetilde{CD}+\widetilde{DA}) & \eta^{6}\equiv\frac{1}{3}( \widetilde{BA}+\widetilde{CB}+\widetilde{DD})\\
     \rho^{7}\equiv\frac{1}{3}( \widetilde{AC}+\widetilde{BB}+\widetilde{CD}) & \eta^{7}\equiv\frac{1}{3}( \widetilde{AB}+\widetilde{BD}+\widetilde{CC})\\
     
     \rho^{8}\equiv\frac{1}{3}( \widetilde{AD}+\widetilde{BA}+\widetilde{CC}) & \eta^{8}\equiv\frac{1}{3}( \widetilde{AA}+\widetilde{BC}+\widetilde{CD})\\
     \rho^{9}\equiv\frac{1}{3}( \widetilde{BA}+\widetilde{CC}+\widetilde{DB}) & \eta^{9}\equiv\frac{1}{3}( \widetilde{BB}+\widetilde{CA}+\widetilde{DC})\\
     \rho^{10}\equiv\frac{1}{3}( \widetilde{AD}+\widetilde{CC}+\widetilde{DB}) & \eta^{10}\equiv\frac{1}{3}( \widetilde{AC}+\widetilde{CB}+\widetilde{DD})\\
     \rho^{11}\equiv\frac{1}{3}( \widetilde{AD}+\widetilde{BA}+\widetilde{DB}) & \eta^{11}\equiv\frac{1}{3}( \widetilde{AB}+\widetilde{BD}+\widetilde{DA})\\
     
     \rho^{12}\equiv\frac{1}{3}( \widetilde{AB}+\widetilde{CA}+\widetilde{DD}) & \eta^{12}\equiv\frac{1}{3}( \widetilde{AA}+\widetilde{CD}+\widetilde{DB})\\
     \rho^{13}\equiv\frac{1}{3}( \widetilde{AB}+\widetilde{BC}+\widetilde{DD}) & \eta^{13}\equiv\frac{1}{3}( \widetilde{AD}+\widetilde{BB}+\widetilde{DC})\\
     \rho^{14}\equiv\frac{1}{3}( \widetilde{AB}+\widetilde{BC}+\widetilde{CA}) & \eta^{14}\equiv\frac{1}{3}( \widetilde{AC}+\widetilde{BA}+\widetilde{CB})\\
     \rho^{15}\equiv\frac{1}{3}( \widetilde{BC}+\widetilde{CA}+\widetilde{DD}) & \eta^{15}\equiv\frac{1}{3}( \widetilde{BD}+\widetilde{CC}+\widetilde{DA})\\

\end{array}
\label{eq: non-commutative quantum numbers}
\end{equation}

The operations associated with the quantum numbers defined in equation (\ref{eq: non-commutative quantum numbers}) are as follows, where the quantum numbers on the left within each set of parentheses is lowered and the one on the right is raised:

\begin{equation}
\begin{array}{ll}
    T^{BD}_{b2}T^{BC}_{b2}:(\rho^0, \eta^0), (\rho^7, \eta^7), (\rho^{10}, \eta^{10}), (\rho^{13}, \eta^{13}) & 
    T^{BC}_{b2}T^{BD}_{b2}:(\eta^0, \rho^0), (\eta^7, \rho^7), (\eta^{10}, \rho^{10}), (\eta^{13}, \rho^{13})\\ 

    T^{CD}_{b2}T^{BD}_{b2}:(\rho^0, \eta^0), (\rho^7, \eta^7), (\rho^{10}, \eta^{10}), (\rho^{13}, \eta^{13}) &
    T^{BD}_{b2}T^{CD}_{b2}:(\eta^0, \rho^0), (\eta^7, \rho^7), (\eta^{10}, \rho^{10}), (\eta^{13}, \rho^{13})\\

    T^{BC}_{b2}T^{CD}_{b2}:(\rho^0, \eta^0), (\rho^7, \eta^7), (\rho^{10}, \eta^{10}), (\rho^{13}, \eta^{13})&
    T^{CD}_{b2}T^{BC}_{b2}:(\eta^0, \rho^0), (\eta^7, \rho^7), (\eta^{10}, \rho^{10}), (\eta^{13}, \rho^{13})\\
    \\
    T^{AC}_{b2}T^{AD}_{b2}:(\rho^2, \eta^2), (\rho^5, \eta^5), (\rho^8, \eta^8), (\rho^{15}, \eta^{15}) &
    T^{AD}_{b2}T^{AC}_{b2}:(\eta^2, \rho^2), (\eta^5, \rho^5), (\eta^8, \rho^8), (\eta^{15}, \rho^{15})\\
    
    T^{AD}_{b2}T^{CD}_{b2}:(\rho^2, \eta^2), (\rho^5, \eta^5), (\rho^8, \eta^8), (\rho^{15}, \eta^{15}) &
    T^{CD}_{b2}T^{AD}_{b2}:(\eta^2, \rho^2), (\eta^5, \rho^5), (\eta^8, \rho^8), (\eta^{15}, \rho^{15})\\
    
    T^{CD}_{b2}T^{AC}_{b2}:(\rho^2, \eta^2), (\rho^5, \eta^5), (\rho^8, \eta^8), (\rho^{15}, \eta^{15}) &
    T^{AC}_{b2}T^{CD}_{b2}:(\eta^2, \rho^2), (\eta^5, \rho^5), (\eta^8, \rho^8), (\eta^{15}, \rho^{15})\\
    \\
    T^{AD}_{b2}T^{AB}_{b2}:(\rho^1, \eta^1), (\rho^6, \eta^6), (\rho^{11}, \eta^{11}), (\rho^{12}, \eta^{12}) &
    T^{AB}_{b2}T^{AD}_{b2}:(\eta^1, \rho^1), (\eta^6, \rho^6), (\eta^{11}, \rho^{11}), (\eta^{12}, \rho^{12})\\

    T^{BD}_{b2}T^{AB}_{b2}:(\rho^1, \eta^1), (\rho^6, \eta^6), (\rho^{11}, \eta^{11}), (\rho^{12}, \eta^{12}) &
    T^{AD}_{b2}T^{BD}_{b2}:(\eta^1, \rho^1), (\eta^6, \rho^6), (\eta^{11}, \rho^{11}), (\eta^{12}, \rho^{12})\\
    
    T^{BD}_{b2}T^{AD}_{b2}:(\rho^1, \eta^1), (\rho^6, \eta^6), (\rho^{11}, \eta^{11}), (\rho^{12}, \eta^{12}) &
    T^{AB}_{b2}T^{BD}_{b2}:(\eta^1, \rho^1), (\eta^6, \rho^6), (\eta^{11}, \rho^{11}), (\eta^{12}, \rho^{12})\\
    \\
    T^{AB}_{b2}T^{AC}_{b2}:(\rho^3, \eta^3), (\rho^4, \eta^4), (\rho^9, \eta^9), (\rho^{14}, \eta^{14}) &
    T^{AC}_{b2}T^{AB}_{b2}:(\eta^3, \rho^3), (\eta^4, \rho^4), (\eta^9, \rho^9), (\eta^{14}, \rho^{14})\\
    
    T^{AC}_{b2}T^{BC}_{b2}:(\rho^3, \eta^3), (\rho^4, \eta^4), (\rho^9, \eta^9), (\rho^{14}, \eta^{14}) &
    T^{BC}_{b2}T^{AC}_{b2}:(\eta^3, \rho^3), (\eta^4, \rho^4), (\eta^9, \rho^9), (\eta^{14}, \rho^{14})\\
    
    T^{BC}_{b2}T^{AB}_{b2}:(\rho^3, \eta^3), (\rho^4, \eta^4), (\rho^9, \eta^9), (\rho^{14}, \eta^{14}) &
    T^{AB}_{b2}T^{BC}_{b2}:(\eta^3, \rho^3), (\eta^4, \rho^4), (\eta^9, \rho^9), (\eta^{14}, \rho^{14})
\end{array}
\end{equation}

The non-local quantum numbers defined in equation (\ref{eq: non-commutative quantum numbers}) are related to those defined in equations (\ref{eq: kappa}) and (\ref{eq: omega}) in the following way:

\begin{equation}
    \begin{array}{ll}
    \kappa^0 = \frac{1}{4}(\rho^0+\rho^1+\rho^2+\rho^3)&
    \kappa^1 = \frac{1}{4}(\rho^4+\rho^5+\rho^6+\rho^7)\\
    \kappa^2 = \frac{1}{4}(\rho^8+\rho^9+\rho^{10}+\rho^{11})&
    \kappa^3 = \frac{1}{4}(\rho^{12}+\rho^{13}+\rho^{14}+\rho^{15})\\
    \\
    \omega^0 = \frac{1}{4}(\eta^0+\eta^4+\eta^8+\eta^{12})&
    \omega^1 = \frac{1}{4}(\eta^1+\eta^5+\eta^9+\eta^{13})\\
    \omega^2 = \frac{1}{4}(\eta^2+\eta^6+\eta^{10}+\eta^{14})&
    \omega^3 = \frac{1}{4}(\eta^3+\eta^7+\eta^{11}+\eta^{15})
    \end{array}
\end{equation}


\section{Chosen non-local operations \label{sec: Chosen non-local operations}}

For the chosen non-local quantum numbers given in equation (\ref{eq: non-local number choice}), the operations which generate odd variations in equation (\ref{eq: interference term}) are as follows:

   \begin{equation}
\begin{array}{ll}
    T^{AB}_{b2}T^{CD}_{b2}: & \; (\nu^0, \mu^0), \; (\kappa^1,\kappa^2), \; (\omega^1,\omega^2)\\
    
    T^{AC}_{b2}T^{BD}_{b2}: & \; (\nu^2, \mu^2), \; (\kappa^1,\kappa^0), \; (\omega^1,\omega^3)\\
    
    T^{AD}_{b2}T^{BC}_{b2}: & \; (\nu^4, \mu^4), \; (\kappa^1,\kappa^3),\;(\omega^1, \omega^0)\\
    
\end{array}
\end{equation}

\begin{equation}
    \begin{array}{ll}
    T^{AC}_{b2}T^{AD}_{b2}:(\rho^2, \eta^2), (\rho^5, \eta^5), (\rho^8, \eta^8), (\rho^{15}, \eta^{15}) &
    T^{AD}_{b2}T^{AC}_{b2}:(\eta^2, \rho^2), (\eta^5, \rho^5), (\eta^8, \rho^8), (\eta^{15}, \rho^{15})\\
    
    T^{AD}_{b2}T^{CD}_{b2}:(\rho^2, \eta^2), (\rho^5, \eta^5), (\rho^8, \eta^8), (\rho^{15}, \eta^{15}) &
    T^{CD}_{b2}T^{AD}_{b2}:(\eta^2, \rho^2), (\eta^5, \rho^5), (\eta^8, \rho^8), (\eta^{15}, \rho^{15})\\
    
    T^{CD}_{b2}T^{AC}_{b2}:(\rho^2, \eta^2), (\rho^5, \eta^5), (\rho^8, \eta^8), (\rho^{15}, \eta^{15}) &
    T^{AC}_{b2}T^{CD}_{b2}:(\eta^2, \rho^2), (\eta^5, \rho^5), (\eta^8, \rho^8), (\eta^{15}, \rho^{15})\\
    \\
    T^{BD}_{b2}T^{BC}_{b2}:(\rho^0, \eta^0), (\rho^{10}, \eta^{10}) & 
    T^{BC}_{b2}T^{BD}_{b2}:(\eta^0, \rho^0), (\eta^{10}, \rho^{10})\\ 

    T^{CD}_{b2}T^{BD}_{b2}:(\rho^0, \eta^0), (\rho^{10}, \eta^{10}) &
    T^{BD}_{b2}T^{CD}_{b2}:(\eta^0, \rho^0), (\eta^{10}, \rho^{10})\\

    T^{BC}_{b2}T^{CD}_{b2}:(\rho^0, \eta^0), (\rho^{10}, \eta^{10})&
    T^{CD}_{b2}T^{BC}_{b2}:(\eta^0, \rho^0), (\eta^{10}, \rho^{10})\\
    \\
    T^{AD}_{b2}T^{AB}_{b2}:(\rho^{11}, \eta^{11}), (\rho^{12}, \eta^{12}) &
    T^{AB}_{b2}T^{AD}_{b2}:(\eta^{11}, \rho^{11}), (\eta^{12}, \rho^{12})\\

    T^{BD}_{b2}T^{AB}_{b2}:(\rho^{11}, \eta^{11}), (\rho^{12}, \eta^{12}) &
    T^{AD}_{b2}T^{BD}_{b2}:(\eta^{11}, \rho^{11}), (\eta^{12}, \rho^{12})\\
    
    T^{BD}_{b2}T^{AD}_{b2}:(\rho^{11}, \eta^{11}), (\rho^{12}, \eta^{12}) &
    T^{AB}_{b2}T^{BD}_{b2}:(\eta^{11}, \rho^{11}), (\eta^{12}, \rho^{12})\\
    \\
    T^{AB}_{b2}T^{AC}_{b2}:(\rho^3, \eta^3), (\rho^{14}, \eta^{14}) &
    T^{AC}_{b2}T^{AB}_{b2}:(\eta^3, \rho^3), (\eta^{14}, \rho^{14})\\
    
    T^{AC}_{b2}T^{BC}_{b2}:(\rho^3, \eta^3), (\rho^{14}, \eta^{14}) &
    T^{BC}_{b2}T^{AC}_{b2}:(\eta^3, \rho^3), (\eta^{14}, \rho^{14})\\
    
    T^{BC}_{b2}T^{AB}_{b2}:(\rho^3, \eta^3), (\rho^{14}, \eta^{14}) &
    T^{AB}_{b2}T^{BC}_{b2}:(\eta^3, \rho^3), (\eta^{14}, \rho^{14})
    \end{array}
\end{equation}


\section{Sample calculations \label{sec: Sample calculations}}
\subsection{Sample calculation: spin systems in rotated frames \label{sec: sample rotated}}

The physical scenario of interest in this sample calculation consists of two rotated Stern-Gerlach detectors which each interact with a particle of total spin $j$. The values of the conditioning variables in this case are $n=6$, $j=1$, $m_{a1}=0$, and $\theta=\frac{\pi}{2}$. In Table \ref{tab: sample epsilons 1}, the cardinalities of Alice's and Bob's elementary state spaces, as defined in equations (\ref{eq: Alice's elementary cardinality}) and (\ref{eq: Bob's elementary cardinality}), are given for all possible combinations of the random variable $m_{b2}$, the local nuisance variables $l_{a1}$ and $l_{b2}$, and the non-local nuisance variable $\nu^0$. In Table \ref{tab: sample local 1}, the post-interference cardinalities of Alice's and Bob's local state spaces, as defined in equations (\ref{eq: local cardinality Alice}) and (\ref{eq: local cardinality Bob}), are given for all possible combinations of the random variable $m_{b2}$ and the local nuisance variables $l_{a1}$ and $l_{b2}$. Also given for each of these combinations are the quantities defined in equations (\ref{eq: Alice local state space configurations}) and (\ref{eq: Bob local state space configurations}), which are the number of unique configurations of Alice's and Bob's local state spaces compatible with fixed $C$ and $D$ symbols in each of their events. The quantity $\Upsilon(n,j,m_{a1},m_{b2},\theta)$, as defined in equation (\ref{eq: upsilon}), can then be found for each possible value of the random variable $m_{b2}$ by summing over the associated products of $|L^a|^\circleddash$, $|L^b|^\circleddash$, $G_{a1}$, and $G_{b2}$:

\begin{equation}
\begin{array}{ll}
    \Upsilon(6,1,0,+1,\frac{\pi}{2})=2592\\
    \Upsilon(6,1,0,0,\frac{\pi}{2})=128\\
    \Upsilon(6,1,0,-1,\frac{\pi}{2})=2592
\end{array}
\end{equation}

We may now calculate the probability $P(m_{b2}|n,j,m_{a1},\theta)$, as defined in equation (\ref{eq: probability}), for each possible value of $m_{b2}$:

\begin{equation}
\begin{array}{ll}
    P(+1|6,1,0,\frac{\pi}{2})=\frac{2592}{2592+128+2592}=0.48795181\\
    P(0|6,1,0,\frac{\pi}{2})=\frac{128}{2592+128+2592}=0.02409639\\
    P(-1|6,1,0,\frac{\pi}{2})=\frac{2592}{2592+128+2592}=0.48795181
\end{array}
\end{equation}

These values can be compared to the predictions of QM, which are $\frac{1}{2}$, $0$, and $\frac{1}{2}$, respectively. As discussed in section \ref{sec: rotated frames results}, the difference between these values will decrease as $n$ increases. 

\subsection{Sample calculation: entangled spin systems \label{sec: sample entangled}}

The physical scenario of interest in this sample calculation consists of two aligned Stern-Gerlach detectors which each interact with a single component of a composite spin system. The values of the conditioning variables in this case are $n=6$, $j_{1a}=1$, $j_{b2}=1$, and $j_{12}=1$, where we enforce the constraint $m_{12}=0$. In Table \ref{tab: sample epsilons 2}, the cardinalities of Alice's and Bob's elementary state spaces, as defined in equations (\ref{eq: Alice's elementary cardinality}) and (\ref{eq: Bob's elementary cardinality}), are given for all possible combinations of the random variables $m_{1a}$ and $m_{b2}$, the local nuisance variable $l_{12}$, and the non-local nuisance variable $\nu^4$. In Table \ref{tab: sample local 2}, the post-interference cardinalities of Alice's and Bob's local state spaces, as defined in equations (\ref{eq: local cardinality Alice}) and (\ref{eq: local cardinality Bob}), are given for all possible combinations of the random variables $m_{1a}$ and $m_{b2}$ and the local nuisance variable $l_{12}$. Also given for each of these combinations are the quantities defined in equations (\ref{eq: Alice local state space configurations}) and (\ref{eq: Bob local state space configurations}), which are the number of unique configurations of Alice's and Bob's local state spaces compatible with fixed $C$ and $D$ symbols in each of their events. The quantity $\Upsilon(n,j_{1a},m_{1a},j_{b2},m_{b2},j_{12})$, as defined in equation (\ref{eq: upsilon}), can then be found for each possible combination of the values of $m_{1a}$ and $m_{b2}$ by summing over the associated products of $|L^a|^\circleddash$, $|L^b|^\circleddash$, $G_{1a}$, and $G_{b2}$:

\begin{equation}
\begin{array}{ll}
    \Upsilon(6,1,+1,1,-1,1)=1280\\
    \Upsilon(6,1,0,1,0,1)=160\\
    \Upsilon(6,1,-1,1,+1,1)=1280
\end{array}
\end{equation}

We may now calculate the probability $P(m_{1a},m_{b2}|n,j_{1a},j_{b2},j_{12})$, as defined in equation (\ref{eq: probability}), for each possible combination of the values of $m_{1a}$ and $m_{b2}$:

\begin{equation}
\begin{array}{ll}
    P(+1,-1|6,1,1,1)=\frac{1280}{1280+160+1280}=0.47058824\\
    P(0,0|6,1,1,1)=\frac{160}{1280+160+1280}=0.05882353\\
    P(-1,+1|6,1,1,1)=\frac{1280}{1280+160+1280}=0.47058824
\end{array}
\end{equation}

These values can be compared to the predictions of QM, which are $\frac{1}{2}$, $0$, and $\frac{1}{2}$, respectively. As discussed in section \ref{sec: entangled results}, the difference between these values will decrease as $n$ increases. 

\subsection{Sample calculation: entangled spin 1/2 systems in rotated frames \label{sec: sample bell}}

The physical scenario of interest in this sample calculation consists of two rotated Stern-Gerlach detectors which each interact with a single component of a composite spin system. The values of the conditioning variables in this case are $n=8$, $j_{1a}=\frac{1}{2}$, $j_{b2}=\frac{1}{2}$, $\theta_{ab}=\frac{\pi}{2}$, and $j_{\gamma}=0$, where we enforce the constraint $m^a_{\gamma}=0$. Due to the large number (64) of unique combinations of quantum numbers involved in this calculation, we do not explicitly list each contribution to $\Upsilon$ as in sections \ref{sec: sample rotated} and \ref{sec: sample entangled}. For each possible combination of the values of $m_{1a}$ and $m_{b2}$, the quantity $\Upsilon(n,j_{1a},m_{1a},j_{b2},m_{b2},\theta_{ab},j_{\gamma})$ is as follows:

\begin{equation}
\begin{array}{ll}
    \Upsilon(8,\frac{1}{2},-\frac{1}{2},\frac{1}{2},-\frac{1}{2},\frac{\pi}{2},0)=51744000\\
    \Upsilon(8,\frac{1}{2},-\frac{1}{2},\frac{1}{2},+\frac{1}{2},\frac{\pi}{2},0)=51744000\\
    \Upsilon(8,\frac{1}{2},+\frac{1}{2},\frac{1}{2},-\frac{1}{2},\frac{\pi}{2},0)=51744000\\
    \Upsilon(8,\frac{1}{2},+\frac{1}{2},\frac{1}{2},+\frac{1}{2},\frac{\pi}{2},0)=51744000
  
\end{array}
\end{equation}

We may now calculate the probability $P(m_{1a},m_{b2}|n,j_{1a},j_{b2},\theta_{ab},j_{\gamma})$, as defined in equation (\ref{eq: probability}), for each possible combination of the values of $m_{1a}$ and $m_{b2}$:

\begin{equation}
\begin{array}{cc}
    P(-\frac{1}{2},-\frac{1}{2}|8,\frac{1}{2},\frac{1}{2},\frac{\pi}{2},0)=\frac{51744000}{51744000+51744000+51744000+51744000}=0.25\\
    P(-\frac{1}{2},+\frac{1}{2}|8,\frac{1}{2},\frac{1}{2},\frac{\pi}{2},0)=\frac{51744000}{51744000+51744000+51744000+51744000}=0.25\\
    P(+\frac{1}{2},-\frac{1}{2}|8,\frac{1}{2},\frac{1}{2},\frac{\pi}{2},0)=\frac{51744000}{51744000+51744000+51744000+51744000}=0.25\\
    P(+\frac{1}{2},+\frac{1}{2}|8,\frac{1}{2},\frac{1}{2},\frac{\pi}{2},0)=\frac{51744000}{51744000+51744000+51744000+51744000}=0.25
\end{array}
\end{equation}

These results agree exactly with the prediction of QM. 

\subsection{Sample calculation: data tables \label{sec: data tables}}

\begin{table}[h]
\begin{centering}
\begin{tabular}{|c|c|c|c|c|c|c|c|c|c|c|c|c|}
\hline 
\multicolumn{13}{|c|}{$n=6$, $j=1$, $m_{a1}=0$, $\theta=\pi/2$}\tabularnewline
\hline 
$m_{b2}$ & $l_{a1}$ & $l_{b2}$ & $\nu^{0}$ & $|\varepsilon^{a}|$ & $|\varepsilon^{b}|$ &  & $m_{b2}$ & $l_{a1}$ & $l_{b2}$ & $\nu^{0}$ & $|\varepsilon^{a}|$ & $|\varepsilon^{b}|$\tabularnewline
\cline{1-6} \cline{8-13} 
-1 & -2 & 0 & 3/4 & 6 & 2 &  & +1 & -2 & 0 & 3/4 & 6 & 2\tabularnewline
-1 & -1 & -1 & 3/4 & 3 & 6 &  & +1 & -1 & -1 & 3/4 & 3 & 6\tabularnewline
-1 & -1 & +1 & 3/4 & 3 & 6 &  & +1 & -1 & +1 & 3/4 & 3 & 6\tabularnewline
-1 & 0 & -2 & 3/4 & 1 & 12 &  & +1 & 0 & -2 & 3/4 & 1 & 12\tabularnewline
-1 & 0 & 0 & 3/4 & 4 & 8 &  & +1 & 0 & 0 & 3/4 & 4 & 8\tabularnewline
-1 & 0 & +2 & 3/4 & 1 & 12 &  & +1 & 0 & +2 & 3/4 & 1 & 12\tabularnewline
-1 & +1 & -1 & 3/4 & 3 & 6 &  & +1 & +1 & -1 & 3/4 & 3 & 6\tabularnewline
-1 & +1 & +1 & 3/4 & 3 & 6 &  & +1 & +1 & +1 & 3/4 & 3 & 6\tabularnewline
-1 & +2 & 0 & 3/4 & 6 & 2 &  & +1 & +2 & 0 & 3/4 & 6 & 2\tabularnewline
0 & -2 & -1 & 5/4 & 4 & 1 &  & 0 & 2 & -1 & 1/4 & 4 & 1\tabularnewline
0 & -2 & +1 & 1/4 & 4 & 1 &  & 0 & 2 & 1 & 5/4 & 4 & 1\tabularnewline
0 & -1 & -2 & 5/4 & 1 & 4 &  & 0 & 1 & -2 & 1/4 & 1 & 4\tabularnewline
0 & -1 & 0 & 5/4 & 3 & 2 &  & 0 & 1 & 0 & 5/4 & 3 & 2\tabularnewline
0 & -1 & 0 & 1/4 & 3 & 2 &  & 0 & 1 & 0 & 1/4 & 3 & 2\tabularnewline
0 & -1 & 2 & 1/4 & 1 & 4 &  & 0 & 1 & 2 & 5/4 & 1 & 4\tabularnewline
0 & 0 & -1 & 5/4 & 2 & 3 &  & 0 & 0 & 1 & 5/4 & 2 & 3\tabularnewline
0 & 0 & -1 & 1/4 & 2 & 3 &  & 0 & 0 & 1 & 1/4 & 2 & 3\tabularnewline
\hline 
\end{tabular}
\par\end{centering}
\caption{The cardinalities of Alice's and Bob's elementary state spaces for all possible combinations of the random variable $m_{b2}$, the local nuisance variables $l_{a1}$ and $l_{b2}$, and the non-local nuisance variable $\nu^0$.}
\label{tab: sample epsilons 1}
\end{table}

\begin{table}[h]
\begin{centering}
\begin{tabular}{|c|c|c|c|c|c|c|c|c|c|c|c|c|c|c|}
\hline 
\multicolumn{15}{|c|}{$n=6$, $j=1$, $m_{a1}=0$, $\theta=\pi/2$}\tabularnewline
\hline 
$m_{b2}$ & $l_{a1}$ & $l_{b2}$ & $|L^{a}|^\circleddash$ & $|L^{b}|^\circleddash$ & $G_{a1}$ & $G_{b2}$ &  & $m_{b2}$ & $l_{a1}$ & $l_{b2}$ & $|L^{a}|^\circleddash$ & $|L^{b}|^\circleddash$ & $G_{a1}$ & $G_{b2}$\tabularnewline
\cline{1-7} \cline{9-15} 
-1 & -2 & 0 & 6 & 2 & 1 & 6 &  & +1 & -2 & 0 & 6 & 2 & 1 & 6\tabularnewline
-1 & -1 & -1 & 3 & 6 & 4 & 4 &  & +1 & -1 & -1 & 3 & 6 & 4 & 4\tabularnewline
-1 & -1 & +1 & 3 & 6 & 4 & 4 &  & +1 & -1 & +1 & 3 & 6 & 4 & 4\tabularnewline
-1 & 0 & -2 & 1 & 12 & 6 & 1 &  & +1 & 0 & -2 & 1 & 12 & 6 & 1\tabularnewline
-1 & 0 & 0 & 4 & 8 & 6 & 6 &  & +1 & 0 & 0 & 4 & 8 & 6 & 6\tabularnewline
-1 & 0 & +2 & 1 & 12 & 6 & 1 &  & +1 & 0 & +2 & 1 & 12 & 6 & 1\tabularnewline
-1 & +1 & -1 & 3 & 6 & 4 & 4 &  & +1 & +1 & -1 & 3 & 6 & 4 & 4\tabularnewline
-1 & +1 & +1 & 3 & 6 & 4 & 4 &  & +1 & +1 & +1 & 3 & 6 & 4 & 4\tabularnewline
-1 & +2 & 0 & 6 & 2 & 1 & 6 &  & +1 & +2 & 0 & 6 & 2 & 1 & 6\tabularnewline
0 & -2 & -1 & 4 & 1 & 1 & 4 &  & 0 & +2 & -1 & 4 & 1 & 1 & 4\tabularnewline
0 & -2 & +1 & 4 & 1 & 1 & 4 &  & 0 & +2 & +1 & 4 & 1 & 1 & 4\tabularnewline
0 & -1 & -2 & 1 & 4 & 4 & 1 &  & 0 & +1 & -2 & 1 & 4 & 4 & 1\tabularnewline
0 & -1 & 0 & 0 & 0 & 4 & 6 &  & 0 & +1 & 0 & 0 & 0 & 4 & 6\tabularnewline
0 & -1 & +2 & 1 & 4 & 4 & 1 &  & 0 & +1 & +2 & 1 & 4 & 4 & 1\tabularnewline
0 & 0 & -1 & 0 & 0 & 6 & 4 &  & 0 & 0 & +1 & 0 & 0 & 6 & 4\tabularnewline
\hline 
\end{tabular}
\par\end{centering}
\caption{The post-interference cardinalities of Alice's and Bob's local state spaces for all possible combinations of the random variable $m_{b2}$ and the local nuisance variables $l_{a1}$ and $l_{b2}$. Also given are the number of unique configurations of these local state spaces compatible with fixed $C$ and $D$ symbols in Alice's and Bob's events.}
\label{tab: sample local 1}
\end{table}

\begin{table}[h]
\begin{centering}
\begin{tabular}{|c|c|c|c|c|c|c|c|c|c|c|c|c|}
\hline 
\multicolumn{13}{|c|}{$n=6$, $j_{1a}=1$, $j_{b2}=1$, $j_{12}=1$, $m_{12}=0$}\tabularnewline
\hline 
$m_{1a}$ & $m_{b2}$ & $l_{12}$ & $\nu^{4}$ & $|\varepsilon^{a}|$ & $|\varepsilon^{b}|$ &  & $m_{1a}$ & $m_{b2}$ & $l_{12}$ & $\nu^{4}$ & $|\varepsilon^{a}|$ & $|\varepsilon^{b}|$\tabularnewline
\cline{1-6} \cline{8-13} 
-1 & +1 & -1 & 1/2 & 8 & 8 &  & +1 & -1 & +1 & 1 & 8 & 8\tabularnewline
-1 & +1 & 0 & 3/4 & 6 & 6 &  & +1 & -1 & 0 & 3/4 & 6 & 6\tabularnewline
-1 & +1 & +1 & 1 & 4 & 4 &  & +1 & -1 & -1 & 1/2 & 4 & 4\tabularnewline
-1 & +1 & +2 & 5/4 & 2 & 2 &  & +1 & -1 & -2 & 1/4 & 2 & 2\tabularnewline
0 & 0 & -2 & 3/4 & 4 & 4 &  & 0 & 0 & +2 & 3/4 & 4 & 4\tabularnewline
0 & 0 & -1 & 0 & 1 & 1 &  & 0 & 0 & +1 & 3/2 & 1 & 1\tabularnewline
0 & 0 & -1 & 1 & 3 & 3 &  & 0 & 0 & +1 & 1/2 & 3 & 3\tabularnewline
0 & 0 & 0 & 1/4 & 2 & 2 &  & 0 & 0 & 0 & 5/4 & 2 & 2\tabularnewline
\hline 
\end{tabular}
\par\end{centering}
\caption{The cardinalities of Alice's and Bob's elementary state spaces for all possible combinations of the random variables $m_{1a}$ and $m_{b2}$, the local nuisance variable $l_{12}$, and the non-local nuisance variable $\nu^4$.}
\label{tab: sample epsilons 2}
\end{table}

\begin{table}[h]
\begin{centering}
\begin{tabular}{|c|c|c|c|c|c|c|c|c|c|c|c|c|c|c|}
\hline 
\multicolumn{15}{|c|}{$n=6$, $j_{1a}=1$, $j_{b2}=1$, $j_{12}=1$, $m_{12}=0$}\tabularnewline
\hline 
$m_{1a}$ & $m_{b2}$ & $l_{12}$ & $|L^{a}|^\circleddash$ & $|L^{b}|^\circleddash$ & $G_{1a}$ & $G_{b2}$ &  & $m_{1a}$ & $m_{b2}$ & $l_{12}$ & $|L^{a}|^\circleddash$ & $|L^{b}|^\circleddash$ & $G_{1a}$ & $G_{b2}$\tabularnewline
\cline{1-7} \cline{9-15} 
-1 & +1 & -1 & 8 & 8 & 1 & 1 &  & +1 & -1 & +1 & 8 & 8 & 1 & 1\tabularnewline
-1 & +1 & 0 & 6 & 6 & 4 & 4 &  & +1 & -1 & 0 & 6 & 6 & 4 & 4\tabularnewline
-1 & +1 & +1 & 4 & 4 & 6 & 6 &  & +1 & -1 & -1 & 4 & 4 & 6 & 6\tabularnewline
-1 & +1 & +2 & 2 & 2 & 4 & 4 &  & +1 & -1 & -2 & 2 & 2 & 4 & 4\tabularnewline
0 & 0 & -2 & 4 & 4 & 1 & 1 &  & 0 & 0 & +2 & 4 & 4 & 1 & 1\tabularnewline
0 & 0 & -1 & 2 & 2 & 4 & 4 &  & 0 & 0 & +1 & 2 & 2 & 4 & 4\tabularnewline
0 & 0 & 0 & 0 & 0 & 6 & 6 &  &  &  &  &  &  &  & \tabularnewline
\hline 
\end{tabular}
\par\end{centering}
\caption{The post-interference cardinalities of Alice's and Bob's local state spaces for all possible combinations of the random variables $m_{1a}$ and $m_{b2}$ and the local nuisance variable $l_{12}$. Also given are the number of unique configurations of these local state spaces compatible with fixed $C$ and $D$ symbols in Alice's and Bob's events.}
\label{tab: sample local 2}
\end{table}


\end{document}